% Vietnamese translation for OI Public Library
% Source-PDF: IOI中国国家候选队论文/2004/韩文弢.pdf
\documentclass[11pt]{article}
\usepackage{oipl-vn}
\usepackage{array}
\usepackage{longtable}

\newcommand{\cpp}{C++}
\newcommand{\pascal}{Pascal}

\title{Bàn về ứng dụng của ngôn ngữ \cpp{} trong thi tin học}
\author{Han Wentao\\Trường Trung học Dư Diêu, Chiết Giang}
\date{IOI 2004, đội tuyển tập huấn quốc gia Trung Quốc}

\begin{document}
\maketitle

\origtitle{Bàn về ứng dụng của ngôn ngữ C++ trong thi tin học}
\sourcepath{IOI China candidate team papers / 2004 / Han Wentao.pdf}

\begin{abstract}
Ngôn ngữ lập trình là một bộ phận quan trọng của thi tin học: mọi thuật toán
chỉ thật sự giải được bài toán sau khi được hiện thực bằng một ngôn ngữ lập
trình. Nhờ tính linh hoạt cao và hệ thống thư viện mạnh, \cpp{} được dùng rất
rộng rãi trong các cuộc thi sinh viên, và cũng ngày càng phổ biến hơn trong
các cuộc thi dành cho học sinh trung học. Bài viết gồm ba chương, lần lượt
giới thiệu từ nông đến sâu những kiến thức cơ bản của \cpp{}, nền tảng lập
trình hướng đối tượng cùng một số thư viện chuẩn thường dùng, và Thư viện
Khuôn mẫu Chuẩn (STL). Hy vọng bài viết có ích cho những độc giả muốn dùng
\cpp{} trong thi tin học.
\end{abstract}

\noindent\textbf{Từ khóa:} thi tin học; ngôn ngữ lập trình \cpp{}; Thư viện
Khuôn mẫu Chuẩn.

\tableofcontents

\section*{Lời nói đầu}
\addcontentsline{toc}{section}{Lời nói đầu}

Thi tin học thường yêu cầu thí sinh trong một thời gian nhất định phải hiểu
và phân tích đề, thiết kế thuật toán thỏa mãn giới hạn thời gian và bộ nhớ,
rồi dùng một ngôn ngữ lập trình để hiện thực thuật toán đúng đắn trên máy
tính. Vì toàn bộ cuộc thi bị giới hạn thời gian, nhất là trong các cuộc thi
kiểu ACM/ICPC, nơi số bài bằng nhau thì thứ hạng phụ thuộc vào tổng thời gian
làm bài, nên ngôn ngữ lập trình có giúp hiện thực thuật toán đúng và nhanh hay
không ảnh hưởng khá lớn đến kết quả. Vì vậy, độ phức tạp khi lập trình ngày
càng được coi trọng; nó lại phụ thuộc rất nhiều vào ngôn ngữ được chọn. Một
số ngôn ngữ thường gặp trong thi tin học có đặc điểm như sau.

\begin{center}
\begin{tabular}{>{\bfseries}lcccc}
\hline
 & BASIC & \pascal{} & \cpp{} & Java\\
\hline
Độ khó học & dễ & trung bình & khó hơn & khó hơn\\
Đặc điểm ngôn ngữ & đơn giản & chặt chẽ & linh hoạt & hướng đối tượng cao\\
Tốc độ chạy & chậm & khá nhanh & nhanh & chậm\\
Chức năng thư viện & yếu & trung bình & rất mạnh & mạnh\\
\hline
\end{tabular}
\end{center}

Trong các cuộc thi tin học học sinh trung học ở thời điểm bài viết, \pascal{}
được dùng khá rộng rãi. Nhưng \cpp{}, nhờ sự linh hoạt vốn có và thư viện rất
mạnh, đang được ngày càng nhiều thí sinh sử dụng. Bài viết này ra đời trong
bối cảnh đó. Trong bài viết, \pascal{} được hiểu theo Free Pascal, vì chuẩn
\pascal{} cung cấp quá ít và Free Pascal phổ biến hơn trong thi đấu; \cpp{}
được hiểu theo chuẩn ANSI/ISO \cpp{}.

Cần chú ý rằng đây không phải là tài liệu tra cứu. Vì vậy, bài viết không thể
giới thiệu mọi mặt của \cpp{} với đầy đủ chi tiết. Muốn học sâu hơn về
\cpp{}, độc giả nên đọc thêm các tài liệu chuyên biệt.

\section{Từ \pascal{} đến \cpp{}}

Điều kiện cần để đọc chương này là đã biết dùng \pascal{} và có một lượng
kiến thức thuật toán nhất định. Vì hiện nay trong thi tin học trung học vẫn có
nhiều thí sinh dùng \pascal{}, mục đích của chương này là giúp độc giả đã nắm
\pascal{} có thể chuyển sang \cpp{} càng nhanh càng tốt. Nếu đã quen với
\cpp{}, có thể bỏ qua chương này và đọc thẳng chương sau.

\subsection{Thế giới, xin chào!}

Trước hết, hãy dùng một chương trình \cpp{} rất đơn giản để có ấn tượng ban
đầu về ngôn ngữ này.

\begin{verbatim}
//: C01:Hello.cpp
// Hello, world!
#include <iostream>
using namespace std;
int main() {
    cout << "Hello, world!" << endl;
} ///:~
\end{verbatim}

Chương trình này in dòng chữ \verb|"Hello, world!"| ra màn hình.

Trong dòng thứ nhất, dòng thứ hai và dòng cuối, phần bắt đầu từ \verb|//| đến
hết dòng đều là chú thích. Trong \cpp{} có hai cách viết chú thích. Một cách
là chú thích khối kế thừa từ C, bắt đầu bằng \verb|/*| và kết thúc bằng
\verb|*/|. Cách còn lại là chú thích dòng, bắt đầu từ \verb|//| và kéo dài
đến cuối dòng.

Nội dung bắt đầu bằng \verb|#| ở dòng thứ ba gọi là chỉ thị tiền xử lý. Dòng
này đưa tập tin đầu \verb|iostream| vào chương trình, tương tự câu lệnh
\verb|uses| để dùng unit trong \pascal{}. Tuy nhiên, \pascal{} mặc định dùng
unit \verb|System|, trong đó chứa phần lớn thủ tục và hàm thường dùng, còn
\cpp{} mặc định không đưa tập tin đầu nào vào. Ngoài ra, tập tin đầu của C
thường có đuôi \verb|.h|, còn \cpp{} chuẩn khuyến khích dùng tập tin đầu
không có phần mở rộng.

Dòng thứ tư cho phép ta trực tiếp dùng các định danh trong không gian tên
\verb|std|. Không gian tên \verb|std| chứa mọi lớp và hàm do \cpp{} chuẩn cung
cấp; để tiện, người ta thường viết dòng này ngay sau các chỉ thị đưa tập tin
đầu vào. Chi tiết về không gian tên có thể xem thêm tài liệu khác.

Dòng thứ năm đến thứ bảy là thân chương trình. Ta định nghĩa một hàm tên
\verb|main|, tương đương chương trình chính trong \pascal{}. Mọi chương trình
\cpp{} có thể chạy độc lập đều cần có hàm \verb|main|. Trong \cpp{}, hai dấu
\verb|{| và \verb|}| tương đương \verb|begin| và \verb|end| của \pascal{}.

Dòng thứ sáu in \verb|Hello, world!| rồi xuống dòng. Ở đây \verb|cout| biểu
diễn đầu ra chuẩn, thường là màn hình, tương tự \verb|Output| trong
\pascal{}. Toán tử \verb|<<| vốn là phép dịch trái bit, nhưng ở đây được định
nghĩa lại thành toán tử chèn, đưa một nội dung vào luồng xuất. \verb|endl| có
nghĩa là xuống dòng. Chuỗi trong \cpp{} được đặt giữa dấu nháy kép, và mỗi
câu lệnh kết thúc bằng dấu chấm phẩy.

\subsection{Kiểu và định nghĩa}

\subsubsection*{Kiểu}

Giống \pascal{}, \cpp{} cung cấp các kiểu cơ bản và cho phép lập trình viên
tự định nghĩa kiểu. Bảng sau liệt kê một số kiểu cơ bản thường dùng trong
\cpp{} cùng kiểu \pascal{} tương ứng.

\begin{center}
\small
\begin{tabular}{lllll}
\hline
\textbf{Tên} & \textbf{Kiểu \cpp{}} & \textbf{Kiểu \pascal{}} &
\textbf{Miền giá trị} & \textbf{Byte}\\
\hline
Logic & \verb|bool| & \verb|Boolean| & \verb|true|/\verb|false| & 1\\
Ký tự & \verb|char| & \verb|Char| & mọi ký tự một byte & 1\\
Số nguyên có dấu 8 bit & \verb|char| & \verb|ShortInt| & \(-128..127\) & 1\\
Số nguyên không dấu 8 bit & \verb|unsigned char| & \verb|Byte| & \(0..255\) & 1\\
Số nguyên có dấu 16 bit & \verb|short| & \verb|SmallInt| & \(-32768..32767\) & 2\\
Số nguyên không dấu 16 bit & \verb|unsigned short| & \verb|Word| & \(0..65535\) & 2\\
Số nguyên có dấu 32 bit & \verb|int| & \verb|LongInt| &
\(-2147483648..2147483647\) & 4\\
Số nguyên không dấu 32 bit & \verb|unsigned int| & \verb|LongWord| &
\(0..4294967295\) & 4\\
Số nguyên có dấu 64 bit & \verb|long long| & \verb|Int64| & \(-2^{63}..2^{63}-1\) & 8\\
Số nguyên không dấu 64 bit & \verb|unsigned long long| & \verb|QWord| & \(0..2^{64}-1\) & 8\\
Dấu phẩy động đơn & \verb|float| & \verb|Single| & \(1.17\mathrm{e}{-38}..3.40\mathrm{e}{38}\) & 4\\
Dấu phẩy động đôi & \verb|double| & \verb|Double| & \(2.22\mathrm{e}{-308}..1.79\mathrm{e}{308}\) & 8\\
Dấu phẩy động mở rộng & \verb|long double| & \verb|Extended| &
\(3.36\mathrm{e}{-4932}..1.18\mathrm{e}{4932}\) & 10/12\\
\hline
\end{tabular}
\end{center}

\paragraph{Kiểu logic.}
Giống \pascal{}, kiểu logic biểu diễn kết quả của các phép toán logic. Một
giá trị kiểu logic chỉ có thể là \verb|true| hoặc \verb|false|; giá trị số
của \verb|true| là \(1\), còn của \verb|false| là \(0\). Khác \pascal{}, trong
\cpp{} nhiều kiểu khác có thể được chuyển ngầm sang kiểu logic: giá trị khác
không trở thành \verb|true|, còn không trở thành \verb|false|.

\paragraph{Kiểu ký tự.}
Hằng ký tự đơn được đặt trong dấu nháy đơn. Một số ký tự không hiển thị có thể
được biểu diễn bằng chuỗi thoát. Ngoài ra, từ bảng trên có thể thấy trong
\cpp{}, kiểu ký tự và kiểu nguyên một byte thực chất tương đương; chẳng hạn
giá trị số của \verb|'A'| là \(65\).

\begin{center}
\begin{tabular}{lll}
\hline
\textbf{Tên} & \textbf{Tên ASCII} & \textbf{Tên trong \cpp{}}\\
\hline
Xuống dòng & NL/LF & \verb|\n|\\
Tab ngang & HT & \verb|\t|\\
Tab dọc & VT & \verb|\v|\\
Lùi một ký tự & BS & \verb|\b|\\
Về đầu dòng & CR & \verb|\r|\\
Sang trang & FF & \verb|\f|\\
Chuông & BEL & \verb|\a|\\
Dấu gạch chéo ngược & \verb|\| & \verb|\\|\\
Dấu hỏi & \verb|?| & \verb|\?|\\
Dấu nháy đơn & \verb|'| & \verb|\'|\\
Dấu nháy kép & \verb|"| & \verb|\"|\\
Mã ASCII bát phân & \verb|ooo| & \verb|\ooo|\\
Mã ASCII thập lục phân & \verb|hhh| & \verb|\xhhh|\\
\hline
\end{tabular}
\end{center}

\paragraph{Kiểu nguyên.}
Kiểu nguyên trong \cpp{} gần giống \pascal{}. Đáng chú ý là tên kiểu nguyên
\cpp{} thực ra gồm ba phần. Thứ nhất là \verb|signed| hoặc \verb|unsigned|,
chỉ có dấu hay không dấu, mặc định là có dấu. Thứ hai là từ bổ nghĩa độ dài
trừ \verb|char|, gồm \verb|short|, \verb|long|, \verb|long long|, mặc định là
\verb|long|. Cuối cùng là \verb|int|, có thể bỏ qua nếu phía trước đã có từ
bổ nghĩa.

Hằng nguyên trong \cpp{} gồm tiền tố, dãy chữ số và hậu tố. Tiền tố mô tả cơ
số: mặc định là thập phân, bắt đầu bằng \verb|0| là bát phân, bắt đầu bằng
\verb|0x| là thập lục phân. Hậu tố mô tả độ dài và tính không dấu:
\verb|l| là số nguyên 32 bit, \verb|ll| là số nguyên 64 bit, \verb|u| là không
dấu, và có thể kết hợp các hậu tố này. Chữ cái trong tiền tố và hậu tố không
phân biệt hoa thường.

\paragraph{Kiểu dấu phẩy động.}
Kiểu dấu phẩy động trong \cpp{} cũng gần giống \pascal{}. Hằng dấu phẩy động
có hai cách viết: dạng cố định và dạng khoa học. Dạng cố định cơ bản giống
\pascal{}, khác ở chỗ trong \cpp{} nếu phần nguyên hoặc phần thập phân bằng
không thì có thể bỏ qua phần đó, nhưng dấu chấm thập phân phải giữ lại; ví dụ
\verb|39.| và \verb|.142857| đều hợp lệ. Dạng khoa học dùng \verb|e| để tách
định trị và số mũ. Hằng dấu phẩy động trong \cpp{} còn có thể mang hậu tố:
\verb|f| biểu diễn kiểu đơn, \verb|l| biểu diễn kiểu mở rộng, không có hậu tố
thì là kiểu đôi.

\paragraph{Kiểu liệt kê.}
\cpp{} cũng có kiểu liệt kê tương tự \pascal{}:

\begin{verbatim}
enum keyword {ASM, AUTO, BREAK}
\end{verbatim}

Kiểu liệt kê này có tên \verb|keyword|, giá trị có thể là \verb|ASM|,
\verb|AUTO| hoặc \verb|BREAK|.

\cpp{} không cung cấp sẵn kiểu miền con và tập hợp. Nó cũng không cung cấp
kiểu chuỗi tích hợp, nhưng có hai cách xử lý chuỗi. Một cách kế thừa từ C,
dựa trên con trỏ ký tự và khá phiền; chi tiết có thể xem tài liệu khác. Cách
còn lại là lớp \verb|string| của thư viện chuẩn \cpp{}, thuận tiện hơn và sẽ
được trình bày ở mục 2.3.

Để biết mỗi kiểu thực tế chiếm bao nhiêu byte, hãy xem chương trình sau.

\begin{verbatim}
//: C01:Types.cpp
// Built-in types
#include <iostream>
using namespace std;
int main() {
    cout << "Size of bool: "
         << sizeof(bool) << endl;
    cout << "Size of char: "
         << sizeof(char) << endl;
    cout << "Size of unsigned char: "
         << sizeof(unsigned char) << endl;
    cout << "Size of short: "
         << sizeof(short) << endl;
    cout << "Size of unsigned short: "
         << sizeof(unsigned short) << endl;
    cout << "Size of int: "
         << sizeof(int) << endl;
    cout << "Size of unsigned int: "
         << sizeof(unsigned int) << endl;
    cout << "Size of long long: "
         << sizeof(long long) << endl;
    cout << "Size of unsigned long long: "
         << sizeof(unsigned long long) << endl;
    cout << "Size of float: "
         << sizeof(float) << endl;
    cout << "Size of double: "
         << sizeof(double) << endl;
    cout << "Size of long double: "
         << sizeof(long double) << endl;
} ///:~
\end{verbatim}

Ở đây \verb|sizeof| trả về kích thước của kiểu trong ngoặc, và cũng có thể áp
dụng cho biến hoặc hằng.

\subsubsection*{Định nghĩa}

Khác \pascal{}, định nghĩa trong \cpp{} gần như có thể xuất hiện ở mọi nơi
trong chương trình. Một định nghĩa gồm bốn phần: từ bổ nghĩa tùy chọn, kiểu cơ
bản, thân định nghĩa và giá trị đầu tùy chọn, thường kết thúc bằng dấu chấm
phẩy. Từ bổ nghĩa thường chỉ định những nội dung không phụ thuộc kiểu. Thân
định nghĩa gồm một tên và một số toán tử bổ nghĩa; các toán tử bổ nghĩa này sẽ
được nói kỹ hơn ở mục 1.3. Ví dụ:

\begin{verbatim}
int number;
char c = 'H';
const double pi = 3.141592653579893;
\end{verbatim}

Ba dòng trên lần lượt định nghĩa biến nguyên \verb|number|, biến ký tự
\verb|c| được khởi tạo bằng \verb|'H'|, và hằng kiểu \verb|double|
\verb|pi| được khởi tạo bằng \(3.141592653579893\). Các biến cùng kiểu có thể
được định nghĩa trong cùng một câu lệnh:

\begin{verbatim}
int a = 39, b, c = 23;
\end{verbatim}

Dòng này định nghĩa ba biến nguyên \verb|a|, \verb|b|, \verb|c|, trong đó
\verb|a| và \verb|c| lần lượt được khởi tạo bằng \(39\) và \(23\).

Với biến không có giá trị đầu, nếu là biến toàn cục, tức được định nghĩa ngoài
mọi hàm, nó sẽ được khởi tạo bằng không theo kiểu tương ứng; nếu là biến cục
bộ, tức được định nghĩa trong một hàm, nó không được tự động khởi tạo.

Tên hay định danh trong \cpp{} có thể gồm chữ cái và chữ số, nhưng ký tự đầu
phải là chữ cái; dấu gạch dưới cũng được xem là một chữ cái. Cần nhớ rằng
khác \pascal{}, \cpp{} phân biệt chữ hoa và chữ thường.

\subsubsection*{Định nghĩa kiểu}

Nếu đặt \verb|typedef| trước một định nghĩa, định nghĩa ấy trở thành định
nghĩa kiểu. Ví dụ:

\begin{verbatim}
typedef unsigned long long QWord;
\end{verbatim}

Khi đó \verb|QWord| không còn là tên biến, mà là tên kiểu. Điều này tương
đương việc làm trong mệnh đề \verb|type| của \pascal{}. Sau đó ta có thể dùng
tên này để định nghĩa biến mới, thường giúp các định nghĩa phía sau ngắn hơn.

\begin{verbatim}
QWord x;
\end{verbatim}

Câu lệnh này định nghĩa biến \verb|x| có kiểu \verb|QWord|, tức
\verb|unsigned long long|.

\subsubsection*{Vòng đời của định nghĩa}

Trong \cpp{}, một tên có thể được dùng sau khi đã được định nghĩa. Tuy nhiên
tên cũng có một vòng đời, hay phạm vi hiệu lực. Khi vượt ra ngoài phạm vi đó,
tên mất hiệu lực và có thể được định nghĩa lại.

Nói chung, vòng đời của biến toàn cục là toàn bộ chương trình. Vòng đời của
biến cục bộ bắt đầu từ chỗ định nghĩa đến khi kết thúc khối chương trình chứa
định nghĩa đó; khối chương trình là một đoạn mã đặt trong cặp ngoặc nhọn. Khi
biến cục bộ trùng tên với biến toàn cục, biến toàn cục bị che khuất và tạm
thời không dùng được bằng tên đó.

\subsection{Con trỏ, mảng và cấu trúc}

\subsubsection*{Con trỏ}

\pascal{} cũng có con trỏ, nhưng con trỏ trong \cpp{} mạnh hơn nhiều. Với mọi
kiểu \(T\), kiểu \verb|T*| biểu diễn kiểu con trỏ trỏ tới một giá trị kiểu
\(T\); dấu \verb|*| ở đây là toán tử bổ nghĩa kiểu đã nói ở mục trước. Ví dụ:

\begin{verbatim}
char c = 'a';
char* p = &c;
\end{verbatim}

Biến \verb|p| là con trỏ tới một giá trị kiểu \verb|char| và được khởi tạo để
trỏ tới biến \verb|c|. Toán tử \verb|&| lấy địa chỉ bộ nhớ của biến phía sau
nó.

Thao tác cơ bản nhất với con trỏ là lấy nội dung mà nó trỏ tới. Trong \cpp{},
thao tác này dùng toán tử \verb|*|.

\begin{verbatim}
char c2 = *p;
\end{verbatim}

Lúc này giá trị của \verb|c2| là \verb|'a'|. Biến con trỏ cũng có thể được gán
bằng \(0\); khi đó nó không trỏ tới ô nhớ hợp lệ nào. Lấy giá trị từ con trỏ
bằng \(0\) sẽ gây lỗi. Con trỏ trong \cpp{} còn hỗ trợ nhiều phép toán khác,
sẽ được mô tả ở mục 1.4.

\subsubsection*{Cấp phát bộ nhớ động}

Giống \pascal{}, \cpp{} có thể cấp phát bộ nhớ động. Việc cấp phát và thu hồi
bộ nhớ được thực hiện bằng toán tử \verb|new| và \verb|delete|, tương tự
\verb|New| và \verb|Dispose| trong \pascal{}, nhưng cách dùng khác nhau.

\begin{verbatim}
int* p = new int;
int* p2 = new int[10];
delete p;
delete[] p2;
\end{verbatim}

Câu lệnh đầu cấp phát không gian cho một biến nguyên và lưu địa chỉ đầu vào
\verb|p|. Câu lệnh thứ hai cấp phát không gian liên tục cho \(10\) biến nguyên
và lưu địa chỉ của biến đầu tiên vào \verb|p2|. Hai câu lệnh sau thu hồi không
gian đã cấp phát. Chú ý ở câu lệnh thứ tư, sau \verb|delete| có cặp ngoặc
vuông rỗng, biểu thị thu hồi không gian của một dãy biến. Khi cấp phát nhiều
biến liên tiếp, số lượng có thể là biến nguyên hoặc hằng nguyên.

\subsubsection*{Con trỏ hằng và con trỏ tới hằng}

Vì con trỏ cũng là một kiểu dữ liệu, nên có biến con trỏ và hằng con trỏ.
Điểm này hơi phức tạp vì ngoài bản thân con trỏ còn có kiểu cơ sở mà nó trỏ
tới. Do đó, hằng liên quan đến con trỏ được chia thành con trỏ hằng và con
trỏ tới hằng.

Con trỏ hằng nghĩa là bản thân con trỏ là hằng: địa chỉ nó trỏ tới không thể
đổi, nhưng nội dung tại địa chỉ đó vẫn có thể đổi. Nói cách khác, \verb|p| là
hằng còn \verb|*p| không phải hằng. Kiểu của con trỏ hằng là \verb|T* const|.

Con trỏ tới hằng nghĩa là bản thân con trỏ không phải hằng: địa chỉ nó trỏ tới
có thể đổi, nhưng nội dung mà nó trỏ tới là hằng và không thể đổi. Nói cách
khác, \verb|p| không phải hằng còn \verb|*p| là hằng. Kiểu này có thể viết
\verb|const T*| hoặc \verb|T const*|.

Hai khái niệm này có thể kết hợp thành con trỏ hằng tới hằng. Khi đó cả
\verb|p| và \verb|*p| đều là hằng, với kiểu \verb|const T* const| hoặc
\verb|T const* const|.

\subsubsection*{Mảng}

Trong \cpp{}, với mọi kiểu \(T\), \verb|T[n]| biểu diễn kiểu mảng gồm \(n\)
phần tử kiểu \(T\). Số \(n\) phải là hằng xác định được lúc biên dịch.

\begin{verbatim}
float v[3];
char* a[32];
\end{verbatim}

Câu lệnh đầu định nghĩa mảng \verb|v| gồm \(3\) phần tử kiểu \verb|float|;
câu lệnh thứ hai định nghĩa mảng \verb|a| gồm \(32\) phần tử kiểu con trỏ ký
tự. Cần nhớ rằng chỉ số mảng trong \cpp{} đều bắt đầu từ \(0\). Vì vậy ba
phần tử của \verb|v| là \verb|v[0]|, \verb|v[1]| và \verb|v[2]|; mảng
\verb|a| cũng tương tự. Điều này có thể làm người quen \pascal{} thấy không
tự nhiên.

Mảng nhiều chiều trong \cpp{} được hiện thực bằng mảng mà phần tử cũng là
mảng.

\begin{verbatim}
int d2[10][20];
\end{verbatim}

Câu lệnh này định nghĩa một mảng hai chiều, chiều thứ nhất kích thước \(10\),
chiều thứ hai kích thước \(20\). Giống \pascal{}, mảng nhiều chiều trong
\cpp{} được lưu theo thứ tự ưu tiên hàng. Đáng chú ý là cả khi định nghĩa và
khi truy cập phần tử, mỗi chiều đều phải đặt trong một cặp ngoặc vuông riêng,
không dùng dấu phẩy để ngăn cách như trong \pascal{}. Thú vị hơn, dấu phẩy
trong \cpp{} cũng là một toán tử: nó lần lượt tính các biểu thức được ngăn
bằng dấu phẩy và trả về giá trị của biểu thức cuối cùng. Vì vậy \verb|d2[1,2]|
không biểu diễn \verb|d2[1][2]|, mà tương đương \verb|d2[2]|, tức một mảng
một chiều gồm \(20\) số nguyên. Độc giả chuyển từ \pascal{} sang \cpp{} cần
tránh lỗi này. Phần lớn trình biên dịch \cpp{} cũng không kiểm tra vượt biên
mảng lúc chạy, nên lập trình viên phải tự chú ý.

\subsubsection*{Khởi tạo mảng}

Khởi tạo mảng tương tự khởi tạo biến thường, nhưng mọi phần tử phải đặt trong
ngoặc nhọn và ngăn cách bằng dấu phẩy.

\begin{verbatim}
int v1[] = {1, 2, 3, 4};
char v2[3] = {'a', 'b', 'c'};
int v3[8] = {1, 2, 3, 4};
\end{verbatim}

Số trong ngoặc vuông có thể bỏ qua. Nếu bỏ qua, kích thước mảng được quyết
định bởi số phần tử trong ngoặc nhọn. Nếu số trong ngoặc vuông lớn hơn số
phần tử khởi tạo, các phần tử cuối còn thiếu giá trị sẽ được khởi tạo bằng
không. Nếu số trong ngoặc vuông nhỏ hơn số phần tử khởi tạo, trình biên dịch
sẽ báo lỗi.

\subsubsection*{Quan hệ giữa mảng và con trỏ}

Trong \cpp{}, con trỏ và mảng có liên hệ rất chặt. Tên của một mảng có thể
được dùng như con trỏ trỏ tới phần tử đầu tiên, hay phần tử thứ không, của
mảng đó.

\begin{verbatim}
int v[] = {1, 2, 3, 4};
int* p1 = v;
int* p2 = &v[0];
int* p3 = &v[4];
\end{verbatim}

Ngược lại, con trỏ cũng có thể được dùng như một mảng.

\begin{verbatim}
int v[] = {1, 2, 3, 4};
int* p = v;
v[2] = 1;
p[2] = 1;
\end{verbatim}

Hai câu lệnh cuối là tương đương.

\subsubsection*{Tham chiếu}

Tham chiếu là bí danh của một biến. Nó chủ yếu được dùng làm tham số hàm và
làm kiểu trả về của một số toán tử được nạp chồng. Với mọi kiểu \(T\),
\verb|T&| biểu diễn kiểu tham chiếu đến \(T\).

\begin{verbatim}
int i = 1;
int& r = i; // r va i bieu dien cung mot bien nguyen
int x = r;  // x = 1
r = 2;      // i = 2
\end{verbatim}

Tham chiếu bắt buộc phải được khởi tạo khi định nghĩa, để bảo đảm nó thật sự
đại diện cho một vùng nhớ xác định.

\subsubsection*{Cấu trúc}

Kiểu cấu trúc trong \cpp{} tương đương kiểu bản ghi trong \pascal{}.

\begin{verbatim}
struct character_type {
    char name[20];
    int level;
    int life;
    int mana;
    int experience;
    int money;
};
\end{verbatim}

Đoạn trên định nghĩa kiểu \verb|character_type|. Giống \pascal{}, biến kiểu
bản ghi có thể dùng dấu chấm để truy cập các thành viên.

\begin{verbatim}
character_type necromancer;
necromancer.level = 99;
necromancer.money = 123456789;
character_type* p = &necromancer;
(*p).life = 1024;
p->mana = 2048;
\end{verbatim}

Với con trỏ trỏ tới một giá trị cấu trúc, ngoài cách \verb|(*ptr).field| ta
cũng có thể dùng \verb|ptr->field|. Không thể viết \verb|*ptr.field|, vì dấu
chấm có độ ưu tiên cao hơn \verb|*|; cách viết đó có nghĩa là \verb|ptr| là
một biến cấu trúc có thành viên \verb|field| kiểu con trỏ, và ta muốn dùng nội
dung mà \verb|ptr.field| trỏ tới.

Trong \cpp{}, định nghĩa danh sách liên kết cũng gọn hơn \pascal{}.

\begin{verbatim}
struct node_type {
    data_type data;
    node_type* next;
};
\end{verbatim}

Tuy vậy, có lúc không tránh khỏi việc các cấu trúc lồng nhau.

\begin{verbatim}
struct list;   // dinh nghia day du se xuat hien sau
struct link {
    list* data;
    link* prev;
    link* next;
};
struct list {
    link* head;
};
\end{verbatim}

Khi đó phải khai báo rỗng trước để đưa tên một kiểu ra trước. Tình huống này
cũng xuất hiện khi các hàm gọi đệ quy lẫn nhau.

\subsection{Biểu thức và câu lệnh}

\subsubsection*{Biểu thức}

Giống \pascal{}, biểu thức trong \cpp{} được tạo bởi toán tử và toán hạng.

\paragraph{Toán hạng.}
Toán hạng trong biểu thức \cpp{} có thể là hằng, hằng số hoặc biến, cũng có
thể là giá trị trả về của hàm. Khác \pascal{}, một số toán tử \cpp{} có tác
dụng phụ, nghĩa là chúng sửa giá trị của toán hạng; các toán tử đó chỉ có thể
tác động lên \term{giá trị trái}. Ban đầu, giá trị trái nghĩa là thứ có thể
đứng bên trái phép gán; đặc điểm của nó là đại diện cho một vùng nhớ có thể
sửa đổi.

\paragraph{Toán tử.}
Bảng sau đối chiếu một số toán tử thường dùng.

\begin{center}
\begin{tabular}{lll}
\hline
\textbf{Tên toán tử} & \textbf{\cpp{}} & \textbf{\pascal{}}\\
\hline
Cộng & \verb|+| & \verb|+|\\
Trừ & \verb|-| & \verb|-|\\
Nhân & \verb|*| & \verb|*|\\
Chia nguyên & \verb|/| & \verb|div|\\
Chia thực & \verb|/| & \verb|/|\\
Lấy dư & \verb|%| & \verb|mod|\\
Nhỏ hơn & \verb|<| & \verb|<|\\
Nhỏ hơn hoặc bằng & \verb|<=| & \verb|<=|\\
Lớn hơn & \verb|>| & \verb|>|\\
Lớn hơn hoặc bằng & \verb|>=| & \verb|>=|\\
Bằng & \verb|==| & \verb|=|\\
Khác & \verb|!=| & \verb|<>|\\
Phủ định bit & \verb|~| & \verb|not|\\
Phủ định logic & \verb|!| & \verb|not|\\
Và bit & \verb|&| & \verb|and|\\
Và logic & \verb|&&| & \verb|and|\\
Hoặc bit & \verb"|" & \verb|or|\\
Hoặc logic & \verb"||" & \verb|or|\\
Xor bit & \verb|^| & \verb|xor|\\
Dịch trái bit & \verb|<<| & \verb|shl|\\
Dịch phải bit & \verb|>>| & \verb|shr|\\
\hline
\end{tabular}
\end{center}

Có thể thấy một đặc điểm lớn của \cpp{} là hầu như mọi toán tử đều được tạo
bởi ký hiệu. Cách dùng các toán tử trong bảng gần giống \pascal{}, nhưng có
ba điểm cần chú ý.

Thứ nhất, trong \cpp{}, chia nguyên và chia thực đều dùng \verb|/|. Khi cả
hai toán hạng là số nguyên thì thực hiện chia nguyên; khi có ít nhất một toán
hạng là số thực thì thực hiện chia thực.

Thứ hai, kiểm tra bằng nhau phải viết hai dấu bằng, vì một dấu bằng được dùng
làm phép gán. Người dùng \pascal{} ban đầu thường không quen điều này. Phiền
hơn nữa, phép gán trong \cpp{} cũng trả về một giá trị. Nếu trong điều kiện
nhầm kiểm tra bằng nhau thành một dấu bằng thì cú pháp vẫn hợp lệ nhưng ý
nghĩa sai; trình biên dịch tốt thường cảnh báo phép gán trong điều kiện.

Thứ ba, toán tử bit và toán tử logic trong \cpp{} là khác nhau.

\cpp{} còn có một số toán tử mà \pascal{} không có. Với toán tử hai ngôi
\verb|@|, biểu thức gán \verb|x = x @ y| có thể viết gọn thành
\verb|x @= y|. Ví dụ \verb|x += 39| tương đương \verb|x = x + 39|. Các toán
tử hỗ trợ kiểu viết này gồm \verb|+|, \verb|-|, \verb|*|, \verb|/|,
\verb|%|, \verb|<<|, \verb|>>|, \verb|&|, \verb"|" và \verb|^|.

\cpp{} cũng có \verb|++| và \verb|--|, đều là toán tử một ngôi, dùng để tăng
hoặc giảm toán hạng một đơn vị. Chúng có dạng tiền tố và hậu tố. Dạng tiền tố
tăng hoặc giảm biến trước rồi trả về giá trị biến; dạng hậu tố trả về giá trị
biến trước rồi mới tăng hoặc giảm. Ví dụ nếu biến nguyên \verb|x| có giá trị
\(12\), biểu thức \verb|++x| có giá trị \(13\), còn \verb|x++| có giá trị
\(12\), nhưng sau đó \verb|x| cũng trở thành \(13\).

Một toán tử ba ngôi thú vị là \verb|?:|. Với biểu thức
\verb|expr1 ? expr2 : expr3|, \cpp{} trước hết xét đúng sai của
\verb|expr1|; nếu đúng thì trả về \verb|expr2|, nếu sai thì trả về
\verb|expr3|.

Chuyển kiểu cũng là một phép toán trong \cpp{}; cả \verb|T(x)| và \verb|(T)x|
đều chuyển \verb|x| sang giá trị kiểu \(T\).

Các toán tử \verb|+|, \verb|-|, \verb|++|, \verb|--|, \verb|+=|, \verb|-=|
có thể dùng với con trỏ. \verb|++| làm con trỏ trỏ tới vị trí kế tiếp ngay
sau ô nhớ hiện tại; xét theo địa chỉ, lượng tăng phụ thuộc vào kích thước của
kiểu mà con trỏ trỏ tới. Các phép toán khác suy ra tương tự. Dấu \verb|-|
còn có thể dùng giữa hai con trỏ cùng kiểu để trả về khoảng cách giữa chúng.
Thao tác chỉ số mảng thực chất là một phép toán con trỏ:
\verb|p[i]| tương đương \verb|*(p + i)|, nên chỉ số cũng là một phép toán
trong \cpp{}. Phép toán con trỏ thể hiện tính linh hoạt của \cpp{}, nhưng
phép toán con trỏ sai có thể khiến chương trình truy cập địa chỉ bị cấm và
sụp đổ.

Về độ ưu tiên toán tử, chỉ cần nhớ một số trường hợp thường gặp, chẳng hạn
nhân chia cao hơn cộng trừ. Khác \pascal{}, trong \cpp{} phép toán logic có độ
ưu tiên thấp hơn phép toán quan hệ, nên một số ngoặc trong điều kiện phức hợp
có thể bỏ qua. Khi không chắc, có thể thêm ngoặc để điều khiển thứ tự tính;
bảng độ ưu tiên đầy đủ nên xem tài liệu tra cứu. Qua hệ thống toán tử có thể
thấy \cpp{} ngắn gọn và linh hoạt hơn \pascal{}.

\subsubsection*{Câu lệnh}

Câu lệnh trong \cpp{} hơi khác \pascal{}. Một câu lệnh có thể là định nghĩa,
một biểu thức đơn, câu lệnh ghép, hoặc câu lệnh điều kiện và lặp sẽ nói dưới
đây. Câu lệnh ghép trong \pascal{} được bao bởi \verb|begin| và \verb|end|;
trong \cpp{} chỉ cần thay bằng ngoặc nhọn trái và phải.

\paragraph{Câu lệnh điều kiện.}
Câu lệnh điều kiện trong \cpp{} cũng là \verb|if|, có hai dạng: có
\verb|else| và không có \verb|else|. Khác \pascal{}, điều kiện cần đặt trong
ngoặc, và sau đó không có \verb|then|.

\begin{verbatim}
if(dieu_kien)
    cau_lenh
\end{verbatim}

hoặc

\begin{verbatim}
if(dieu_kien)
    cau_lenh
else
    cau_lenh
\end{verbatim}

\cpp{} còn có câu lệnh \verb|switch| tương ứng với \verb|case| trong
\pascal{}.

\begin{verbatim}
switch(bieu_thuc) {
    case hang:
        cau_lenh
        break;
    case hang:
        cau_lenh
        break;
    ...
    default:
        cau_lenh
}
\end{verbatim}

Sau mỗi \verb|case| có thể có nhiều câu lệnh đơn, và \verb|break| cuối cùng
không bắt buộc. Nếu sau một \verb|case| không có \verb|break|, sau khi thực
hiện các câu lệnh của \verb|case| đó chương trình sẽ tiếp tục thực hiện
\verb|case| ngay sau, cho đến khi gặp \verb|break| hoặc hết \verb|switch|.
Điều này tiện khi muốn xử lý nhiều giá trị giống nhau.

\begin{verbatim}
switch(c) {
    case '0': case '1': case '2': case '3': case '4':
    case '5': case '6': case '7': case '8': case '9':
        // xu ly chu so
        break;
    ...
}
\end{verbatim}

Nếu muốn mỗi \verb|case| kết thúc là thoát khỏi \verb|switch|, phải thêm
\verb|break|.

\paragraph{Câu lệnh lặp.}
Giống \pascal{}, \cpp{} có ba loại vòng lặp, nhưng cú pháp khác nhiều.

\begin{verbatim}
while(dieu_kien)
    cau_lenh
\end{verbatim}

\begin{verbatim}
do
    cau_lenh
while(dieu_kien)
\end{verbatim}

Hai vòng lặp này đều tiếp tục chạy khi điều kiện đúng và dừng khi điều kiện
sai. Khác biệt là vòng thứ nhất kiểm tra trước mỗi lượt lặp, còn vòng thứ hai
kiểm tra sau mỗi lượt lặp.

Vòng \verb|for| là một trong những câu lệnh khác biệt nhất giữa \cpp{} và
\pascal{}, đồng thời thể hiện rõ tính linh hoạt của \cpp{}.

\begin{verbatim}
for(khoi_tao; dieu_kien; thay_doi)
    cau_lenh
\end{verbatim}

Các phần trong ngoặc được ngăn cách bằng dấu chấm phẩy. Trước khi vòng lặp
bắt đầu, câu lệnh \verb|for| thực hiện phần khởi tạo, thường dùng để khởi tạo
biến lặp; biến được định nghĩa trong phần khởi tạo có vòng đời đến hết vòng
lặp. Trước mỗi lượt lặp, \verb|for| kiểm tra điều kiện; nếu đúng thì thực
hiện thân vòng lặp, nếu sai thì thoát. Sau mỗi lượt lặp, phần thay đổi được
thực hiện, thường dùng để sửa biến lặp.

\begin{verbatim}
int sum = 0;
for(int i = 1; i <= 100; i++)
    sum += i;
\end{verbatim}

Tương tự \pascal{}, cả ba loại vòng lặp đều có thể dùng \verb|break| và
\verb|continue| để thay đổi luồng điều khiển.

\paragraph{Khi nào dùng dấu chấm phẩy.}
Nói chung, mỗi câu lệnh trong \cpp{} cần dấu chấm phẩy để kết thúc, kể cả câu
lệnh nằm giữa \verb|if| và \verb|else|; điểm này khác \pascal{}. Nhưng nếu đó
là câu lệnh ghép, sau ngoặc nhọn phải không cần dấu chấm phẩy. Điều này chỉ
áp dụng cho ngoặc nhọn biểu diễn câu lệnh ghép; nếu là định nghĩa kiểu liệt
kê, định nghĩa cấu trúc hoặc khởi tạo mảng, sau ngoặc nhọn ở vị trí vốn cần
dấu chấm phẩy vẫn phải có dấu chấm phẩy.

\subsection{Hàm}

Trong \cpp{}, định nghĩa một hàm có dạng:

\begin{verbatim}
kieu_tra_ve ten_ham(danh_sach_tham_so) {
    cau_lenh
}
\end{verbatim}

So với \pascal{}, \cpp{} không có khái niệm thủ tục; thủ tục chính là hàm có
kiểu trả về \verb|void|, tức không trả về gì. Khác \pascal{}, hàm \cpp{} không
thể được định nghĩa lồng trong một hàm khác. Khi vào một hàm, biến có thể dùng
chỉ gồm biến toàn cục và các biến chỉ ra trong danh sách tham số. Trong
\pascal{}, nếu thủ tục hoặc hàm không có tham số thì bỏ ngoặc của danh sách
tham số, còn trong \cpp{} ngoặc luôn luôn phải có. Các mục trong danh sách
tham số ngăn cách bằng dấu phẩy, hình thức giống định nghĩa biến.

\begin{verbatim}
int f(int a, int b) {
    ...
}
\end{verbatim}

Tương tự thao tác chỉ số, lời gọi hàm trong \cpp{} cũng được xem là một phép
toán; các toán hạng là bản thân hàm và các tham số thực, kết quả là giá trị
trả về của hàm.

\subsubsection*{Cách truyền tham số}

Trong \pascal{} có tham trị và tham biến. Trong \cpp{}, tham số thông thường
được truyền theo kiểu tham trị, còn tham số định nghĩa bằng tham chiếu được
truyền theo kiểu tham biến. Vì truyền tham trị phải sao chép toàn bộ giá trị
của biến, với biến lớn mà hàm bảo đảm không sửa giá trị, ta có thể truyền
bằng tham chiếu hằng.

\begin{verbatim}
int f(int a, int& b, const int& c) {
    ...
}
\end{verbatim}

Ở đây \verb|a| được truyền theo tham trị, \verb|b| được truyền theo tham biến,
còn \verb|c| cũng truyền địa chỉ biến nhưng nhờ \verb|const| nên trong hàm
\verb|f| giá trị của \verb|c| không thể bị thay đổi. Mọi câu lệnh trong
\verb|f| cố sửa \verb|c| sẽ khiến trình biên dịch báo lỗi. Giống \pascal{},
khi gọi hàm, vị trí tham biến chỉ có thể thay bằng biến, hay chính xác hơn là
giá trị trái. Tuy nhiên, nếu tham số là tham chiếu hằng thì biểu thức thông
thường cũng được; hệ thống sẽ tự tạo biến tạm cho nó.

Mảng trong \cpp{} được truyền dưới dạng con trỏ. Vì vậy mọi thay đổi một phần
tử mảng trong thân hàm sẽ làm thay đổi phần tử tương ứng của mảng lúc gọi,
vì chúng có cùng địa chỉ trong bộ nhớ.

\subsubsection*{Trả về giá trị hàm}

Khác \pascal{}, giá trị hàm trong \cpp{} được trả bằng câu lệnh
\verb|return|. Câu lệnh này vừa trả về giá trị hàm, vừa kết thúc việc thực
hiện hàm và quay lại vị trí gọi hàm. Dạng của nó là:

\begin{verbatim}
return gia_tri_tra_ve;
\end{verbatim}

Hàm kiểu \verb|void| có thể dùng trực tiếp \verb|return;| để kết thúc.

\subsubsection*{Hàm nội tuyến}

Thêm \verb|inline| trước định nghĩa hàm thường sẽ tạo ra định nghĩa hàm nội
tuyến. Khi đó trình biên dịch cố gắng chép trực tiếp các lệnh trong hàm vào
nơi gọi, thay vì truyền tham số và nhảy lệnh. Vì vậy, các hàm đơn giản và
được dùng nhiều lần thường được định nghĩa nội tuyến.

\begin{verbatim}
inline int max(int x, int y) {
    if(x > y)
        return x;
    else
        return y;
}
\end{verbatim}

Nếu hàm nội tuyến không thể được mở trực tiếp, chẳng hạn trong hàm có cấu
trúc phức tạp như vòng lặp, hiệu quả sẽ giống hàm thường.

\subsubsection*{Nạp chồng hàm và tham số mặc định}

\cpp{} hỗ trợ nạp chồng hàm và tham số mặc định; cả hai đều thể hiện sự ngắn
gọn và linh hoạt của ngôn ngữ.

Nạp chồng hàm nghĩa là các hàm cùng tên có thể được định nghĩa nhiều lần,
nhưng mỗi lần có danh sách tham số khác nhau. Khi hàm được gọi, trình biên
dịch dựa vào danh sách tham số thực để quyết định gọi hàm nào. Thông thường
các hàm được nạp chồng dùng để làm cùng một việc, chỉ khác đôi chút về tham
số.

\begin{verbatim}
void print(int i) {
    cout << i << endl;
}
void print(const char* s) {
    cout << s << endl;
}
int main() {
    print(39);             // goi ham thu nhat
    print("Hello, world!"); // goi ham thu hai
}
\end{verbatim}

Tham số mặc định là việc chỉ định trước giá trị mặc định cho một số tham số
khi định nghĩa hàm. Khi gọi, nếu các tham số này có giá trị bằng mặc định thì
có thể bỏ qua. Các tham số có mặc định chỉ được đặt ở cuối danh sách tham số.

\begin{verbatim}
void inc(int& x, int d = 1) {
    x += d;
}
int main() {
    int i = 1;
    inc(i);    // tuong duong inc(i, 1);
    inc(i, 3);
}
\end{verbatim}

Với cả nạp chồng hàm và tham số mặc định, cần tránh tạo ra sự mơ hồ. Vì hai
tính năng này không quá quan trọng trong thi tin học, bài viết chỉ nhắc qua;
chi tiết nên xem tài liệu khác.

\subsection{Một số hàm thư viện thường dùng}

Như đã nói trong lời nói đầu, \cpp{} chuẩn cung cấp thư viện rất mạnh. Mục
này chỉ giới thiệu vài hàm thư viện có chức năng gần với các thủ tục và hàm
chuẩn của \pascal{}.

\begin{center}
\scriptsize
\begin{tabular}{p{0.33\textwidth}p{0.12\textwidth}p{0.31\textwidth}p{0.15\textwidth}}
\hline
\textbf{Định nghĩa hàm} & \textbf{Tập tin đầu} & \textbf{Tác dụng} & \textbf{Ghi chú}\\
\hline
\verb|void* memset(void* p, int b, size_t n);| & \verb|cstring| &
Đặt \(n\) byte liên tiếp bắt đầu tại \verb|p| thành giá trị \verb|b| &
Giống \verb|FillChar|, chú ý thứ tự tham số\\
\verb|void* memmove(void* p, const void* q, size_t n);| & \verb|cstring| &
Sao chép \(n\) byte từ vùng \verb|q| đến vùng \verb|p| &
Giống \verb|Move|; hai vùng có thể chồng lấn một phần\\
\verb|double atof(const char* p);|, \verb|int atoi(const char* p);|,
\verb|long atol(const char* p);| & \verb|cstdlib| &
Chuyển chuỗi \verb|p| thành số được biểu diễn & Giống \verb|Val|\\
\verb|double fabs(double);| & \verb|cmath| & Hàm trị tuyệt đối & Giống \verb|Abs|\\
\verb|double ceil(double);|, \verb|double floor(double);| & \verb|cmath| &
Lấy nguyên, lần lượt là làm tròn lên và làm tròn xuống & \\
\verb|double sqrt(double);| & \verb|cmath| & Căn bậc hai & Giống \verb|Sqrt|\\
\verb|double pow(double d, double e);| & \verb|cmath| & Lũy thừa, trả về \(d^e\) & \\
\verb|double sin(double);|, \verb|double cos(double);|, \verb|double tan(double);| &
\verb|cmath| & Hàm lượng giác & \\
\verb|double asin(double);|, \verb|double acos(double);|, \verb|double atan(double);| &
\verb|cmath| & Hàm lượng giác ngược & \\
\verb|double atan2(double y, double x);| & \verb|cmath| &
Hàm arctan tăng cường, trả về góc cực của điểm \((x,y)\) &
Rất hữu ích, tự xét góc phần tư\\
\verb|double sinh(double);|, \verb|double cosh(double);|, \verb|double tanh(double);| &
\verb|cmath| & Hàm hyperbol & \\
\verb|double exp(double);| & \verb|cmath| & Hàm mũ cơ số \(e\) & Giống \verb|Exp|\\
\verb|double log(double);|, \verb|double log10(double);| & \verb|cmath| &
Logarit cơ số \(e\) và cơ số \(10\) & \verb|log| giống \verb|Ln|\\
\hline
\end{tabular}
\end{center}

\cpp{} chuẩn không cung cấp hàm \verb|Pi|. Muốn lấy giá trị \(\pi\), thường
viết:

\begin{verbatim}
const double pi = acos(0.) * 2;
\end{verbatim}

Chi tiết hơn về các hàm thư viện nên xem tài liệu liên quan.

\subsection{Kết chương}

Sau khi nắm các mục trên, về cơ bản ta đã có thể dùng \cpp{} để làm những
việc thường làm bằng \pascal{}. Mục này dùng \cpp{} để giải một bài thực tế:
bài 3 bảng nâng cao NOIP 2003, \emph{Cây nhị phân cộng điểm}.

\paragraph{Mô tả bài toán.}
Cho một cây nhị phân \(T\) có \(n\) đỉnh, duyệt giữa của nó là
\((1,2,3,\ldots,n)\), trong đó các số là mã đỉnh. Mỗi đỉnh có một điểm là số
nguyên dương; điểm của đỉnh thứ \(j\) ký hiệu \(d_j\). Cây \(T\) và mọi cây
con của nó đều có một giá trị cộng điểm. Với cây con \(S\), giá trị này bằng
giá trị cộng điểm của cây con trái nhân với giá trị cộng điểm của cây con
phải, rồi cộng điểm của gốc \(S\).

Nếu cây con rỗng thì giá trị cộng điểm quy ước là \(1\). Với lá, giá trị cộng
điểm chính là điểm của lá, không xét hai cây con rỗng. Hãy tìm một cây nhị
phân có duyệt giữa \((1,2,3,\ldots,n)\) và giá trị cộng điểm lớn nhất; cần in
giá trị lớn nhất và duyệt trước của cây đó.

\paragraph{Dữ liệu vào.}
Dòng thứ nhất là số nguyên \(n<30\), số đỉnh. Dòng thứ hai gồm \(n\) số nguyên
cách nhau bởi khoảng trắng, là điểm của từng đỉnh, mỗi điểm nhỏ hơn \(100\).

\paragraph{Dữ liệu ra.}
Dòng thứ nhất là giá trị cộng điểm lớn nhất, kết quả không vượt quá
\(4\,000\,000\,000\). Dòng thứ hai gồm \(n\) số nguyên cách nhau bởi khoảng
trắng, là duyệt trước của cây.

\begin{verbatim}
Input
5
5 7 1 2 10

Output
145
3 1 2 4 5
\end{verbatim}

Để đơn giản, chương trình dưới đây dùng bàn phím và màn hình. Nhập chuẩn
trong \cpp{} cũng rất đơn giản: \verb|cin| biểu diễn đầu vào chuẩn, thường là
bàn phím, và thao tác trích bằng \verb|>>| thực hiện nhập.

Từ mô tả bài toán có thể thấy khi tìm giá trị cộng điểm lớn nhất của \(T\),
có nhiều bài toán con là tìm giá trị của các cây con, và các bài toán con
này cũng thỏa mãn tính tối ưu. Trong duyệt giữa, các đỉnh của cùng một cây
con có chỉ số liên tiếp, nên chỉ cần chỉ ra vị trí đầu và cuối là mô tả đầy
đủ một bài toán con. Vì vậy có thể dùng quy hoạch động. Gọi \(\mathrm{best}_{i,j}\)
là giá trị cộng điểm lớn nhất của cây con từ đỉnh \(i\) đến đỉnh \(j\). Khi
đó:
\[
\mathrm{best}_{i,j}=
\begin{cases}
d_j, & 1\le i=j\le n,\\
\max\left\{
\mathrm{best}_{i+1,j}+d_i,\,
\mathrm{best}_{i,j-1}+d_j,\,
\max_{i<k<j}\left(\mathrm{best}_{i,k-1}\mathrm{best}_{k+1,j}+d_k\right)
\right\}, & 1\le i<j\le n.
\end{cases}
\]

Vì đề còn yêu cầu in duyệt trước của cây có điểm lớn nhất, ta dùng
\(\mathrm{root}_{i,j}\) để ghi chỉ số gốc của cây con từ \(i\) đến \(j\). Cuối
cùng dùng đệ quy để in duyệt trước.

\begin{verbatim}
//: C01:Tree.cpp
// Binary Tree with Scores (NOIP2003 Advanced Division)
#include <iostream>
#include <cstring>
using namespace std;
const int max_nodes = 30;
int n;
unsigned int best[max_nodes][max_nodes];
int root[max_nodes][max_nodes];
bool first;
void visit(int i, int j) {
    if(i > j)
        return;
    if(first)
        first = false;
    else
        cout << ' ';
    cout << root[i][j] + 1;
    visit(i, root[i][j] - 1);
    visit(root[i][j] + 1, j);
}
int main() {
    cin >> n;
    memset(best, 0, sizeof(best));
    memset(root, 0xFF, sizeof(root));
    for(int i = 0; i < n; i++) {
        cin >> best[i][i];
        root[i][i] = i;
    }
    for(int d = 1; d < n; d++)
        for(int i = 0; i + d < n; i++) {
            int j = i + d;
            if(best[i + 1][j] > best[i][j - 1]) {
                best[i][j] = best[i][i] + best[i + 1][j];
                root[i][j] = i;
            }
            else {
                best[i][j] = best[j][j] + best[i][j - 1];
                root[i][j] = j;
            }
            for(int k = i + 1; k < j; k++) {
                unsigned int value = best[k][k]
                    + best[i][k - 1] * best[k + 1][j];
                if(value > best[i][j]) {
                    best[i][j] = value;
                    root[i][j] = k;
                }
            }
        }
    cout << best[0][n - 1] << endl;
    first = true;
    visit(0, n - 1);
    cout << endl;
}
///:~
\end{verbatim}

Cần chú ý rằng chỉ số mảng trong \cpp{} bắt đầu từ \(0\), nên chương trình đã
điều chỉnh chỉ số tương ứng. Qua chương này có thể thấy \cpp{} và \pascal{}
có nhiều điểm giống nhau, nhưng cũng có nhiều khác biệt. Ta nên dùng phương
pháp loại suy: lấy điểm giống làm nền tảng, rồi nắm chắc các điểm khác biệt;
thông qua luyện tập liên tục, cuối cùng sẽ sử dụng \cpp{} thành thạo.

\section{Đi sâu hơn vào ngôn ngữ \cpp{}}

Điều kiện cần để đọc chương này là biết viết chương trình \cpp{} đơn giản và
có một lượng kiến thức thuật toán nhất định. Chương này giới thiệu lớp và đối
tượng trong \cpp{}, cùng một số lớp thường dùng của thư viện chuẩn.

\subsection{Lớp}

Lớp không được dùng trực tiếp quá nhiều trong thi tin học. Tuy nhiên, các thư
viện mạnh của \cpp{} đều được xây dựng trên cơ sở lớp, nên ta cần có cái nhìn
cơ bản về chúng. Trong mục này, hai thuật ngữ \term{lớp} và \term{đối tượng}
xuất hiện thường xuyên. Lớp thực chất là một kiểu, còn đối tượng là một biến
có kiểu là một lớp nào đó.

\subsubsection*{Định nghĩa lớp}

Nói đơn giản, lớp là sự mở rộng của kiểu cấu trúc. Thêm một số hàm vào cấu
trúc là có lớp đơn giản nhất.

\begin{verbatim}
struct Date {
    int d, m, y;
    void init(int dd, int mm, int yy);
    void add_year(int n);
    void add_month(int n);
    void add_day(int n);
};
\end{verbatim}

Khi đó ta có thể truy cập hàm thành viên giống như truy cập biến thành viên.

\begin{verbatim}
Date national_day;
national_day.init(1, 10, 2003);
Date next_day;
next_day = national_day;
next_day.add_day(1);
\end{verbatim}

Khi định nghĩa hàm thành viên của lớp ở bên ngoài thân lớp, để tránh trùng tên
hàm giữa các lớp khác nhau, phải thêm tên lớp trước tên hàm và ngăn cách bằng
\verb|::|.

\begin{verbatim}
void Date::init(int dd, int mm, int yy) {
    d = dd;
    m = mm;
    y = yy;
}
\end{verbatim}

Trong hàm thành viên của lớp, muốn dùng biến thành viên chỉ cần viết trực
tiếp tên biến.

\subsubsection*{Quyền truy cập}

Một tư tưởng quan trọng của lập trình hướng đối tượng là đóng gói dữ liệu,
tức giấu dữ liệu bên trong đối tượng và chỉ thao tác lên đối tượng thông qua
hàm thành viên. Khi đó cần chỉ định quyền truy cập cho các thành viên trong
lớp.

\begin{verbatim}
struct Date {
private:
    int d, m, y;
public:
    void init(int dd, int mm, int yy);
    void add_year(int n);
    void add_month(int n);
    void add_day(int n);
};
\end{verbatim}

Ở đây, các biến thành viên \verb|d|, \verb|m|, \verb|y| có quyền riêng tư,
chỉ hàm thành viên của chính lớp mới truy cập được. Bốn hàm thành viên phía
dưới có quyền công khai, bên ngoài đều có thể gọi. Dĩ nhiên không phải biến
thành viên lúc nào cũng phải riêng tư, hay hàm thành viên lúc nào cũng phải
công khai; một số hàm thành viên dùng nội bộ cũng có thể là riêng tư, và biến
thành viên cũng có thể công khai.

Khi định nghĩa lớp, ta thường dùng \verb|class| thay cho \verb|struct|.

\begin{verbatim}
class Date {
private:
    int d, m, y;
public:
    void init(int dd, int mm, int yy);
    void add_year(int n);
    void add_month(int n);
    void add_day(int n);
};
\end{verbatim}

Khác biệt duy nhất giữa \verb|class| và \verb|struct| là khi không chỉ định
quyền truy cập, \verb|struct| mặc định công khai còn \verb|class| mặc định
riêng tư.

\subsubsection*{Hàm dựng và hàm hủy}

Trong một lớp có thể định nghĩa hai loại hàm đặc biệt. Loại thứ nhất là hàm
dựng: không có kiểu trả về, cùng tên với lớp, được gọi tự động khi đối tượng
được tạo, thường dùng để khởi tạo đối tượng; sau danh sách tham số của nó có
thể có danh sách khởi tạo. Loại thứ hai là hàm hủy: cũng không có kiểu trả
về, tên gồm dấu \verb|~| đặt trước tên lớp, được gọi tự động khi đối tượng kết
thúc vòng đời, thường dùng để dọn dẹp. Hàm dựng có thể có tham số và có thể
được nạp chồng; hàm hủy không có tham số.

\begin{verbatim}
class Date {
public:
    Date(int dd, int mm, int yy);
    Date(char* date);
    ~Date();
};
Date::Date(int dd, int mm, int yy) : d(dd), m(mm), y(yy) {
}
Date national_day(1, 10, 2003);
Date christmas("December 25th, 2003");
\end{verbatim}

Trong hiện thực hàm dựng của lớp \verb|Date|, phần sau dấu hai chấm là danh
sách khởi tạo, có thể dùng để khởi tạo biến thành viên. Cũng có thể khởi tạo
trong thân hàm dựng, nhưng thông thường nên dùng danh sách khởi tạo.

\subsubsection*{Thành viên tĩnh}

Trong một lớp, có những biến thành viên được chia sẻ bởi mọi đối tượng của
lớp. Có thể dùng biến toàn cục, nhưng làm như vậy phá vỡ đóng gói. Để giải
quyết, ta định nghĩa chúng là biến thành viên tĩnh. Cũng có những hàm thành
viên có thể được gọi mà không gắn với một đối tượng cụ thể; ta định nghĩa
chúng là hàm thành viên tĩnh. Chỉ cần thêm \verb|static| trước định nghĩa
thành viên thông thường.

Thành viên tĩnh vẫn tuân theo quyền truy cập của lớp. Biến thành viên tĩnh
phải được khởi tạo bên ngoài định nghĩa lớp. Hàm thành viên tĩnh không được
dùng biến thành viên thường, cũng không được gọi hàm thành viên thường, nhưng
có thể dùng biến thành viên tĩnh và gọi hàm thành viên tĩnh khác. Hàm thành
viên thường có thể dùng biến thành viên tĩnh và gọi hàm thành viên tĩnh. Ví
dụ sau minh họa cách định nghĩa và dùng thành viên tĩnh.

\begin{verbatim}
//: C02:Counter.cpp
// Static members in class
#include <iostream>
using namespace std;
class Counter {
public:
    Counter(int id2);
    ~Counter();
    int get_id();
    static int get_tot();
private:
    int id;
    static int tot;
};
int Counter::tot = 0;
Counter::Counter(int id2) :id(id2) {
    tot++;
}
Counter::~Counter() {
    tot--;
}
int Counter::get_id() {
    return id;
}
int Counter::get_tot() {
    return tot;
}
int main() {
    Counter c1(0), c2(1);
    cout << c1.get_id() << endl;
    cout << c1.get_tot() << endl;
    Counter c3(2);
    cout << Counter::get_tot() << endl;
} ///:~
\end{verbatim}

\subsubsection*{Hàm thành viên hằng}

Trong một lớp, có những hàm thành viên không thay đổi trạng thái đối tượng,
chẳng hạn \verb|get_id()| của lớp \verb|Counter| ở trên. Khi đó có thể thêm
\verb|const| sau danh sách tham số của định nghĩa hàm thành viên để khai báo
chúng là hàm thành viên hằng. Trong hàm thành viên hằng, không được viết câu
lệnh sửa biến thành viên; nếu có, trình biên dịch sẽ báo lỗi.

Ngoài ra, đối tượng hằng chỉ có thể dùng hàm thành viên hằng. Ví dụ ở chương
trước, nếu truyền đối tượng bằng tham chiếu hằng, trong hàm chỉ có thể dùng
các hàm thành viên hằng của đối tượng đó. Hàm thành viên tĩnh không thể được
định nghĩa là hàm thành viên hằng. Sửa chương trình trên như sau:

\begin{verbatim}
//: C02:Counter2.cpp
// Constant member functions in class
#include <iostream>
using namespace std;
class Counter {
public:
    Counter(int id2) : id(id2) {
        tot++;
    }
    ~Counter() {
        tot--;
    }
    int get_id() const {
        return id;
    };
    static int get_tot() {
        return tot;
    }
private:
    int id;
    static int tot;
};
int Counter::tot = 0;
int main() {
    const Counter c1(0), c2(1);
    cout << c1.get_id() << endl;
    cout << c1.get_tot() << endl;
    Counter c3(2);
    cout << Counter::get_tot() << endl;
} ///:~
\end{verbatim}

Từ chương trình này cũng thấy rằng hiện thực hàm thành viên có thể đặt ngay
trong định nghĩa lớp. Nếu làm vậy, trình biên dịch sẽ cố xử lý chúng theo
cách nội tuyến. Thông thường các hàm thành viên ngắn được hiện thực đầy đủ
trong định nghĩa lớp, còn hàm dài thì hiện thực bên ngoài.

\subsubsection*{Con trỏ this}

Đôi khi trong hàm thành viên của một đối tượng ta cần tham chiếu đến chính
đối tượng đó; khi ấy dùng con trỏ \verb|this|. Trong hàm thành viên,
\verb|this| là con trỏ tới kiểu của đối tượng và luôn trỏ tới đối tượng hiện
tại. Ngoài ra, nếu trong hàm thành viên có biến cục bộ trùng tên với biến
thành viên, viết trực tiếp tên biến sẽ chỉ biến cục bộ. Muốn truy cập biến
thành viên thì dùng \verb|this|, chẳng hạn \verb|this->x|.

Trên đây là phần giới thiệu ngắn về lớp trong \cpp{}. Thực ra lợi ích lớn
nhất của lớp và đối tượng là đơn giản hóa đáng kể việc viết và bảo trì phần
mềm lớn. Trong tin học, dùng lớp và đối tượng đúng mức cũng có thể giảm độ
phức tạp khi lập trình.

\subsection{Nạp chồng toán tử}

Nạp chồng toán tử nghĩa là định nghĩa lại các toán tử mà ngôn ngữ cung cấp.
Mục đích ban đầu của kỹ thuật này chỉ là làm mã ngắn gọn hơn. Chẳng hạn ta
viết một lớp biểu diễn số nguyên độ chính xác cao, và đối tượng của lớp này
cần hỗ trợ nhiều phép toán số học. Nếu mọi phép toán đều viết thành hàm, mã
sẽ kém đẹp: biểu thức đơn giản \verb|x + y * z| có thể phải viết thành
\verb|add(x, multiply(y, z))|. Nếu có thể định nghĩa lại toán tử của ngôn
ngữ, ta dùng kiểu tự định nghĩa giống như dùng kiểu dựng sẵn.

Như chương trước đã nói, toán tử trong \cpp{} rất phong phú, đến cả chỉ số và
lời gọi hàm cũng có thể xem là toán tử. Khi nạp chồng toán tử, ngoài toán tử
số học thông thường, hầu như mọi toán tử đều có thể được định nghĩa lại.

Dĩ nhiên nạp chồng toán tử có một số ràng buộc. Trước hết, chỉ có thể nạp
chồng toán tử đã tồn tại trong \cpp{}, không thể tự định nghĩa toán tử mới.
Thứ hai, toán tử được nạp chồng phải có ít nhất một toán hạng thuộc kiểu do
người dùng định nghĩa; tức không thể nạp chồng toán tử mà mọi toán hạng đều
là kiểu dựng sẵn. Thứ ba, toán tử sau khi nạp chồng vẫn giữ độ ưu tiên ban
đầu. Ngoài ra, một số toán tử như \verb|.| và \verb|?:| không thể nạp chồng.

Nói chung, chức năng của toán tử được nạp chồng nên phù hợp với ý nghĩa ban
đầu của nó. Điều này không tuyệt đối: trong nhập xuất bằng luồng, toán tử
trích \verb|>>| và toán tử chèn \verb|<<| đã thay đổi hoàn toàn chức năng ban
đầu, nhưng sự thay đổi ấy rất khéo.

Có thể xem toán tử là một loại hàm đặc biệt, còn toán hạng là tham số của hàm.
Vì vậy ta có thể nạp chồng toán tử giống như định nghĩa hàm. Tên hàm của toán
tử là \verb|operator| cộng với chính toán tử đó. Với toán tử hai ngôi
\verb|@|, biểu thức \verb|a @ b| tương đương \verb|operator @(a, b)| hoặc
\verb|a.operator @(b)|. Với toán tử một ngôi \verb|@|, biểu thức \verb|@a|
tương đương \verb|operator @(a)| hoặc \verb|a.operator @()|. Nói cách khác,
toán tử được nạp chồng có thể là hàm toàn cục hoặc hàm thành viên của lớp.

\begin{verbatim}
class complex {
    double re, im;
public:
    complex(double r, double i) : re(r), im(i) {}
    double real() const { return re; }
    double imag() const { return im; }
    complex operator +(const complex& x) {
        return complex(re + x.re, im + x.im);
    }
};
complex operator -(const complex& x, const complex& y) {
    return complex(x.real() - y.real(), x.imag() - y.imag());
}
\end{verbatim}

Trong đoạn trên, toán tử cộng được định nghĩa là hàm thành viên, còn toán tử
trừ được định nghĩa là hàm toàn cục. Thông thường toán tử hai ngôi phổ thông
được định nghĩa thành hàm toàn cục, toán tử một ngôi được định nghĩa thành
hàm thành viên. Tuy nhiên các toán tử trả về kết quả là giá trị trái, như
toán tử gán, thường cũng được định nghĩa thành hàm thành viên và kiểu trả về
phải là tham chiếu.

\subsubsection*{Một số toán tử đặc biệt}

\paragraph{Toán tử chỉ số.}
Toán tử chỉ số \verb|[]| bắt buộc phải được định nghĩa là hàm thành viên của
lớp, thường được nạp chồng khi định nghĩa lớp giống mảng. Điều thú vị là sau
khi nạp chồng, giá trị trong ngoặc vuông không nhất thiết là số nguyên; điều
này rất thuận tiện cho lớp ánh xạ trong STL.

\paragraph{Toán tử lời gọi hàm.}
Toán tử lời gọi hàm \verb|()| bắt buộc phải được định nghĩa là hàm thành viên
của lớp; số lượng và kiểu tham số không bị hạn chế. Nói chung, nó thường được
dùng để định nghĩa mảng nhiều chiều hoặc các đối tượng có thể sử dụng như
hàm.

\paragraph{Toán tử tăng giảm.}
Toán tử \verb|++| và \verb|--| có hai cách dùng: tiền tố và hậu tố. Nạp chồng
dạng tiền tố giống toán tử một ngôi thông thường. Dạng hậu tố được thêm một
tham số nguyên giả, nhờ đó có thể phân biệt với dạng tiền tố theo quy tắc nạp
chồng hàm.

\begin{verbatim}
class counter {
    int v;
public:
    counter(int val = 0) : v(val) { }
    int get_value() const { return v; }

    counter& operator ++() {   // dang tien to
        v++;
        return *this;
    }
    counter operator ++(int) { // dang hau to
        counter temp = *this;
        v++;
        return temp;
    }
};
\end{verbatim}

Các ví dụ cụ thể về nạp chồng toán tử sẽ còn xuất hiện ở những phần sau.

\subsection{Chuỗi}

Thư viện chuẩn \cpp{} cung cấp lớp \verb|string| để xử lý chuỗi thông thường.
Lớp này được định nghĩa trong tập tin đầu \verb|string|, vì vậy khi dùng đừng
quên đưa tập tin đầu này vào.

Thực ra để hỗ trợ các bộ ký tự khác, \cpp{} trước hết dùng khuôn mẫu để hiện
thực một lớp chuỗi tổng quát \verb|basic_string|; \verb|string| chỉ là một bản
đặc biệt của nó với kiểu ký tự là \verb|char|. Để tiện, mục này xem
\verb|string| như một lớp độc lập và giới thiệu một số hàm thường dùng.

Trước hết là vài điểm cơ bản. Chuỗi của lớp \verb|string| về lý thuyết không
bị giới hạn độ dài; nó tự động cấp phát không gian theo độ dài thực tế.
Kiểu \verb|string::size_type| thực chất là \verb|unsigned int|, dùng để mô tả
vị trí và độ dài trong chuỗi. \verb|string::npos| là hằng thành viên tĩnh có
kiểu \verb|size_type|, giá trị là giá trị lớn nhất của \verb|unsigned int|,
thường dùng để mô tả độ dài lớn nhất. Với đối tượng \verb|string| \verb|s|,
ta dùng \verb|s[i]| để truy cập ký tự thứ \(i\); giống chỉ số mảng trong
\cpp{}, ký tự cũng được đánh số từ \(0\).

Nhờ nạp chồng toán tử, phép gán, nối và so sánh của \verb|string| không khác
mấy kiểu \verb|string| trong \pascal{}. Đáng chú ý là nơi cần một chuỗi khác
không chỉ nhận đối tượng \verb|string|, mà cả con trỏ ký tự kế thừa từ C cũng
có thể dùng. Ngoài ra, hàm thành viên \verb|c_str()| trả về một con trỏ ký tự
có nội dung giống đối tượng \verb|string|. Các thủ tục hoặc hàm chuỗi của
\pascal{} đều có hàm thành viên tương ứng trong \verb|string|.

\begin{center}
\begin{tabular}{lll}
\hline
\textbf{\pascal{}} & \textbf{Tác dụng} & \textbf{\texttt{string}}\\
\hline
\verb|Length| & Lấy độ dài chuỗi & \verb|length|\\
\verb|Copy| & Lấy chuỗi con & \verb|substr|\\
\verb|Insert| & Chèn một đoạn chuỗi & \verb|insert|\\
\verb|Delete| & Xóa một đoạn chuỗi & \verb|erase|\\
\verb|Pos| & Tìm vị trí xuất hiện của chuỗi con & \verb|find|\\
\hline
\end{tabular}
\end{center}

Các hàm thành viên này đều có nhiều phiên bản nạp chồng để dùng linh hoạt.
Dưới đây là các cách dùng thường gặp nhất:

\begin{verbatim}
size_type length() const;
string substr(size_type pos = 0, size_type n = npos) const;
string& insert(size_type pos, const string& s);
string& erase(size_type pos = 0, size_type n = npos);
size_type find(const string& s) const;
\end{verbatim}

Trong danh sách tham số trên, \verb|pos| là vị trí trong chuỗi gốc,
\verb|n| là độ dài chuỗi con, còn \verb|s| là một chuỗi khác. Nếu \verb|n|
lớn hơn độ dài còn lại của chuỗi, chuỗi con sẽ đi từ \verb|pos| đến hết chuỗi
gốc. Nếu không chỉ định \verb|n|, tham số mặc định \verb|npos| sẽ lấy chuỗi
con đến hết chuỗi gốc.

Khác \pascal{}, vì \(0\) là vị trí hợp lệ trong chuỗi \verb|string|, khi tìm
chuỗi con mà không thấy, \verb|find| trả về \verb|npos| chứ không trả về
\(0\). Hàm \verb|rfind| cũng tìm chuỗi con, nhưng khác \verb|find| ở chỗ nó
trả về vị trí xuất hiện cuối cùng, còn \verb|find| trả về vị trí xuất hiện
đầu tiên. Trên đây chỉ là các hàm thành viên thường dùng nhất của
\verb|string|; chi tiết cần xem tài liệu.

\subsection{Luồng}

Trong \cpp{}, nhập xuất được hiện thực bằng luồng. Thiết kế các lớp luồng của
thư viện chuẩn \cpp{} thể hiện rõ tính linh hoạt của ngôn ngữ và ưu thế của
lập trình hướng đối tượng.

Trong \cpp{}, lớp cơ sở của mọi luồng nhập là \verb|istream|, lớp cơ sở của
mọi luồng xuất là \verb|ostream|. Luồng nhập chuẩn \verb|cin|, luồng nhập tập
tin và luồng nhập chuỗi đều dẫn xuất từ \verb|istream|, nên nơi nào dùng được
\verb|istream| thì các luồng nhập này đều dùng được; luồng xuất cũng tương tự.

Thao tác thường dùng nhất với luồng là nhập và xuất. Trong các ví dụ trước,
với một luồng nhập ta dùng toán tử trích \verb|>>| để nhập, còn với một luồng
xuất ta dùng toán tử chèn \verb|<<| để xuất. Các thao tác này gần như dùng
được với mọi kiểu dựng sẵn của ngôn ngữ. Cần chú ý rằng chúng tự động xử lý
định dạng dữ liệu theo kiểu. Đặc biệt khi nhập ký tự và chuỗi, gồm cả con trỏ
ký tự và lớp \verb|string|, toán tử trích trước hết bỏ qua mọi ký tự trắng
như khoảng trắng, tab và xuống dòng, rồi mới trích dữ liệu. Với chuỗi, khi gặp
lại ký tự trắng thì thao tác trích kết thúc.

Với một số yêu cầu nhập đặc biệt, ta có thể dùng nhập không định dạng. Các
hàm thành viên thường dùng liên quan đến nhập không định dạng gồm:

\begin{verbatim}
int get();
istream& get(int& c);
istream& get(char* p, int n, char term = '\n');
istream& getline(char* p, int n, char term = '\n');
\end{verbatim}

Hai hàm đầu đều đọc một ký tự từ luồng, kể cả ký tự trắng. Chú ý chỗ lẽ ra
là \verb|char| lại dùng \verb|int|, chủ yếu vì khi đọc đến cuối luồng chúng
phải trả về \(-1\). Hai hàm sau đọc nhiều ký tự từ luồng, tối đa \(n\) ký tự,
nhưng nếu gặp \verb|term| thì kết thúc sớm; cả hai đòi hỏi vùng \verb|p| có
đủ không gian. Khác biệt là \verb|get()| vẫn để ký tự \verb|term| trong luồng,
còn \verb|getline()| bỏ luôn ký tự đó, điều này ảnh hưởng đến các lần nhập
sau.

Hàm thành viên \verb|peek()| trả về ký tự kế tiếp sẽ được đọc từ luồng nhưng
không thật sự lấy nó ra, tức ký tự đó vẫn nằm trong luồng. Hàm \verb|unget()|
đẩy ký tự vừa đọc gần nhất trở lại luồng. Hàm \verb|eof()| kiểm tra luồng đã
đến cuối chưa, giống \verb|Eof| của \pascal{}. Có lúc cũng có thể dùng trực
tiếp đối tượng luồng. Ở nơi cần giá trị logic, chẳng hạn trong ngoặc của
\verb|if|, \cpp{} tự chuyển luồng thành giá trị logic biểu thị luồng còn bình
thường không. Nếu đã đọc đến cuối luồng, giá trị logic này thành sai. Ví dụ,
đoạn sau liên tục đọc số nguyên đến cuối luồng:

\begin{verbatim}
int i;
while(cin >> i) {
    // Do something
}
\end{verbatim}

Các hàm thành viên của luồng không có hàm giống \verb|Eoln| trong \pascal{},
nhưng có thể kiểm tra cuối dòng bằng \verb|cin.peek() == '\n'|.

\subsubsection*{Xuất định dạng}

Trong \pascal{}, khi xuất ta có thể chỉ định độ rộng trường và số chữ số phần
thập phân. Luồng của \cpp{} cũng làm được. Chỉ định độ rộng trường dùng hàm
thành viên \verb|width()|. Chỉ định số chữ số phần thập phân hơi phiền hơn:
trước hết đặt cách xuất số thực về dạng cố định, sau đó đặt số chữ số phần
thập phân.

\begin{verbatim}
cout.setf(ios::fixed, ios::floatfield);
cout.precision(2);
cout << 1.2345 << endl;
\end{verbatim}

Câu lệnh đầu đặt cách xuất số thực thành dạng cố định, câu lệnh thứ hai đặt
số chữ số phần thập phân là \(2\). Giống \pascal{}, chữ số cuối của phần thập
phân được làm tròn.

Cần chú ý \verb|width()| chỉ có hiệu lực với lần xuất định dạng kế tiếp. Nếu
nhiều lần xuất cần chỉ định độ rộng trường thì phải gọi \verb|width()| nhiều
lần. Còn \verb|precision()| có hiệu lực với mọi lần xuất số thực sau đó.

\begin{verbatim}
cout.width(3);
cout << 1 << 2; // 1 co do rong 3, 2 dung do rong thuc
cout.width(4);
cout << 2;
\end{verbatim}

Nếu thấy cách viết này phiền, có thể viết:

\begin{verbatim}
cout << setw(3) << 1 << 2 << setw(4) << 2;
\end{verbatim}

Khi đó đừng quên thêm \verb|#include <iomanip>| ở đầu chương trình, vì
\verb|setw| được định nghĩa trong \verb|iomanip|. Các nội dung xuất định dạng
của luồng \cpp{} còn nhiều, có thể xem thêm tài liệu.

\subsubsection*{Luồng tập tin và luồng chuỗi}

Một ưu điểm lớn của việc \cpp{} trừu tượng hóa nhập xuất thành luồng là ta
không cần quan tâm luồng đang thao tác tồn tại dưới hình thức nào. Khi nhập,
nguồn của luồng là bàn phím, tập tin hay nội dung trong một chuỗi thì ngoài
cách tạo luồng ban đầu có chút khác biệt, các thao tác sau đó đều giống nhau.

\paragraph{Luồng tập tin.}
Trong \cpp{}, lớp của luồng nhập tập tin là \verb|ifstream|, lớp của luồng
xuất tập tin là \verb|ofstream|, cả hai được định nghĩa trong tập tin đầu
\verb|fstream|.

Để mở một tập tin, có thể chỉ định tên tập tin khi tạo luồng, hoặc tạo xong
rồi dùng hàm thành viên \verb|open()|.

\begin{verbatim}
ifstream fin("test.in"); // chi dinh khi tao, dung duoc ngay
ofstream fout;
fout.open("test.out");   // chi dinh sau khi tao
\end{verbatim}

Sau khi luồng tập tin được mở, ta dùng chúng giống như \verb|cin| hoặc
\verb|cout|. Thông thường không cần viết riêng câu lệnh đóng luồng tập tin,
vì hàm hủy của lớp luồng tập tin sẽ tự động đóng tập tin. Nếu muốn đóng trước
khi hàm hủy chạy, có thể dùng hàm thành viên \verb|close()|.

\paragraph{Luồng chuỗi.}
Luồng chuỗi trong \cpp{} dùng để đọc dữ liệu từ một chuỗi, hoặc xuất dữ liệu
vào một chuỗi. Lớp của luồng nhập chuỗi là \verb|istringstream|, lớp của
luồng xuất chuỗi là \verb|ostringstream|, cả hai được định nghĩa trong tập tin
đầu \verb|sstream|.

Khi tạo luồng nhập chuỗi, ta cần chỉ định một chuỗi làm nguồn dữ liệu. Trong
quá trình sử dụng, có thể gọi hàm thành viên \verb|str()| bất cứ lúc nào để
lấy chuỗi tương ứng. Luồng chuỗi có thể dùng để chuyển đổi giữa dữ liệu kiểu
cụ thể và chuỗi.

\begin{verbatim}
string s("123");
istringstream sin(s);
int i;
sin >> i;
cout << i << endl;
ostringstream sout;
sout << i;
string s2 = sout.str();
cout << s2 << endl;
\end{verbatim}

\subsubsection*{Nhập xuất kiểu tự định nghĩa}

Ta đã biết luồng \cpp{} cung cấp toán tử trích và chèn cho các kiểu dữ liệu
dựng sẵn, giúp nhập xuất thuận tiện. Với kiểu tự định nghĩa, ta cũng có thể
dùng \verb|>>| và \verb|<<| để nhập xuất hay không? Câu trả lời là có, và rất
đơn giản: nạp chồng \verb|>>| và \verb|<<|. Ví dụ:

\begin{verbatim}
//: C02:Complex.cpp
// Input & Output for User-defined types
#include <iostream>
using namespace std;
class Complex {
public:
    Complex(double r = 0., double i = 0.)
        : re(r), im(i) {}
    double real() const { return re; }
    double imag() const { return im; }
private:
    double re, im;
    friend istream& operator >>(
        istream& is, Complex& c);
    friend ostream& operator <<(
        ostream& os, const Complex& c);
};
istream& operator >>(istream& is, Complex& c) {
    is >> ws;
    if(is.peek() == '(') {
        char chr;
        is >> chr;    // bo qua '('
        is >> c.re;
        is >> chr;    // bo qua ','
        is >> c.im;
        is >> chr;    // bo qua ')'
    }
    else {
        is >> c.re;
        c.im = 0.;
    }
    return is;
}
ostream& operator <<(ostream& os, const Complex& c) {
    return os << '(' << c.re << ',' << c.im << ')';
}
int main() {
    Complex c;
    cin >> c;
    cout << c << endl;
} ///:~
\end{verbatim}

Chương trình trên hiện thực một kiểu số phức đơn giản và nạp chồng toán tử
trích, chèn. Vì ta không thể sửa mã của lớp luồng, hai toán tử này chỉ có thể
được nạp chồng toàn cục. Kiểu trả về là tham chiếu đến luồng, nhờ đó các toán
tử trích hoặc chèn có thể dùng liên tiếp. Trong danh sách tham số, tham số đầu
tiên phải là tham chiếu đến luồng. Khi nạp chồng toán tử trích, tham số thứ
hai cũng phải là tham chiếu, vì khi nhập ta cần thay đổi giá trị biến. Khi
nạp chồng toán tử chèn, tham số thứ hai thường là tham chiếu hằng, vừa bảo
đảm tốc độ vì tham chiếu truyền như con trỏ, vừa bảo đảm biến không bị sửa.
Khi kết thúc hàm, đừng quên trả về chính luồng để toán tử có thể dùng liên
tiếp.

\subsection{Kết chương}

Có lẽ độc giả đã thấy ở đầu và cuối các chương trình trong bài có vài chú
thích lạ. Chúng dùng để trích mã nguồn của từng chương trình từ tập tin văn
bản thuần của bài viết. Sau đây là một chương trình thực hiện chức năng đó.

\begin{verbatim}
//: C02:Extract.cpp
// Extracts codes from documents
#include <iostream>
#include <fstream>
#include <string>
#include <cstdlib>
using namespace std;
void error(const char* msg) {
    cerr << msg << endl;
    exit(1); // similar to Pascal's Halt
}
int main(int argc, char* argv[]) {
    if(argc <= 1)
        error("Document filename missing");
    ifstream fin(argv[1]);
    if(!fin)
        error("Cannot open file");
    string line;
    bool is_in_code = false;
    ofstream fout;
    while(getline(fin, line))
        if(line.find("//:") == 0) {
            if(is_in_code) {
                fout.close();
                error("Begin tag unexpected");
            }
            else {
                is_in_code = true;
                string name = line.substr(line.rfind(':') + 1);
                fout.open(name.c_str());
                if(!fout)
                    error("Cannot create file");
                fout << line << endl;
            }
        }
        else if(line.find("///:~") != string::npos) {
            if(is_in_code) {
                is_in_code = false;
                fout << line << endl;
                fout.close();
            }
            else
                error("End tag unexpected");
        }
        else if(is_in_code)
            fout << line << endl;
    if(is_in_code) {
        fout.close();
        error("End tag expected");
    }
    fin.close();
} ///:~
\end{verbatim}

Chương trình này minh họa chuỗi và luồng đã giới thiệu trong chương. Cách
dùng rất đơn giản: sau khi biên dịch, gõ tên chương trình trên dòng lệnh và
thêm tên tập tin văn bản thuần của tài liệu làm tham số, chương trình sẽ tự
động trích các chương trình nguồn trong bài ra thư mục hiện tại.

Trong chương này, ta đã giới thiệu cơ chế trừu tượng hóa dữ liệu của \cpp{},
tức lớp, rồi từ đó giới thiệu một số đặc tính ngôn ngữ xây dựng trên nền lớp.
Chúng đều là công cụ mạnh cho việc lập trình.

\section{Giới thiệu STL}

Điều kiện cần để đọc chương này là hiểu nền tảng lập trình hướng đối tượng
của \cpp{} và có một lượng kiến thức thuật toán nhất định. Chương này giới
thiệu STL, tức \emph{Standard Template Library}, Thư viện Khuôn mẫu Chuẩn
mạnh mẽ do \cpp{} chuẩn cung cấp.

\subsection{Khái quát STL}

Đôi khi ta thấy một số hàm có cùng chức năng nhưng tham số khác kiểu, khiến
ta phải viết nhiều hàm tương tự, rất bất tiện. Để giải quyết vấn đề này,
\cpp{} cung cấp phương tiện lập trình tổng quát, trong đó kiểu cũng được xem
như tham số. Ví dụ hàm thường dùng để đổi giá trị hai biến:

\begin{verbatim}
void swap(int& x, int& y) {
    int t = x;
    x = y;
    y = t;
}
\end{verbatim}

Hàm trên chỉ đổi được hai biến nguyên. Với kiểu khác, ta buộc phải viết hàm
tương tự. Vậy dùng lập trình tổng quát để đơn giản hóa ra sao?

\subsubsection*{Hàm khuôn mẫu}

Trong \cpp{}, vấn đề này được giải quyết bằng hàm khuôn mẫu. Khi dùng hàm
khuôn mẫu, ta chỉ cần sửa hàm cũ rất ít. Với \verb|swap|, có thể viết:

\begin{verbatim}
template <class T> void swap(T& x, T& y) {
    T t = x;
    x = y;
    y = t;
}
\end{verbatim}

Trong hàm \verb|swap| trên, \verb|template| là tiền tố của hàm khuôn mẫu,
biểu thị phần sau là một hàm khuôn mẫu. Nội dung giữa \verb|<| và \verb|>|
mô tả tham số khuôn mẫu; \verb|class T| nghĩa là \(T\) là một kiểu, không nhất
thiết là lớp. Trong định nghĩa hàm, ta dùng \(T\) thay cho \verb|int| ban đầu.
Sau khi định nghĩa hàm khuôn mẫu \verb|swap|, ta có thể dùng nó để đổi giá
trị của mọi cặp biến cùng kiểu; trình biên dịch sẽ dựa vào kiểu khác nhau để
tạo các phiên bản \verb|swap| khác nhau, qua đó giảm đáng kể độ phức tạp khi
lập trình.

Thông thường trình biên dịch dựa vào kiểu tham số thực khi gọi để đặc biệt
hóa hàm khuôn mẫu. Ta cũng có thể tự chỉ định kiểu tham số khi đặc biệt hóa.

\begin{verbatim}
swap<int>(x, y);
\end{verbatim}

Ở đây \verb|swap<int>()| thực chất giống hàm ban đầu đổi hai biến nguyên.

\subsubsection*{Lớp khuôn mẫu}

Tương tự, ta cũng có thể tổng quát hóa lớp.

\begin{verbatim}
template <class T, int max> struct c_array {
    typedef T value_type;
    typedef T& reference;
    typedef const T& const_reference;
    T v[max];
    operator T*() { return v; }
    reference operator [](size_t i) { return v[i]; }
    const_reference operator [](size_t i) const { return v[i]; }
    size_t size() const { return max; }
};
\end{verbatim}

Đoạn trên định nghĩa lớp khuôn mẫu \verb|c_array|, có thể thay cho mảng một
chiều dựng sẵn của \cpp{}. Nó có hai tham số khuôn mẫu: tham số thứ nhất chỉ
kiểu phần tử, tham số thứ hai chỉ số phần tử. Nếu muốn định nghĩa một mảng
nguyên gồm \(100\) phần tử, viết:

\begin{verbatim}
c_array<int, 100> a;
\end{verbatim}

Vì đã nạp chồng toán tử chỉ số và toán tử chuyển kiểu, ta có thể dùng lớp này
giống mảng thông thường.

\subsubsection*{Thư viện Khuôn mẫu Chuẩn}

STL là thư viện mạnh được xây dựng trên hàm khuôn mẫu và lớp khuôn mẫu. Nhờ
hàm khuôn mẫu, ta có thể viết các hàm hỗ trợ gần như mọi kiểu, kể cả kiểu tự
định nghĩa, để hiện thực các thuật toán tổng quát thường dùng như đếm, sắp
xếp, tìm kiếm. Nhờ lớp khuôn mẫu, ta có thể viết các vật chứa hỗ trợ gần như
mọi kiểu, dùng để hiện thực các cấu trúc dữ liệu thường gặp như danh sách
liên kết, ngăn xếp, hàng đợi, cây nhị phân cân bằng.

Điều thú vị là khác các thư viện thông thường, các tập tin đầu tạo nên STL
chứa toàn bộ mã nguồn của các hàm và lớp. Ta có thể đọc các tập tin đầu này
để học cách STL hiện thực các chức năng. Các mục sau sẽ lần lượt giới thiệu
các thành phần của STL.

\subsection{Bộ lặp}

Bộ lặp thực chất là một kiểu con trỏ tổng quát, tức sự trừu tượng hóa của
kiểu con trỏ. Dựa trên các thao tác được hỗ trợ, bộ lặp được chia thành năm
loại: bộ lặp xuất, bộ lặp nhập, bộ lặp tiến, bộ lặp hai chiều và bộ lặp ngẫu
nhiên. Bảng sau liệt kê các thao tác được hỗ trợ.

\begin{center}
\scriptsize
\begin{tabular}{llllll}
\hline
\textbf{Loại bộ lặp} & \textbf{Xuất} & \textbf{Nhập} & \textbf{Tiến} &
\textbf{Hai chiều} & \textbf{Ngẫu nhiên}\\
\hline
Viết tắt & Out & In & For & Bi & Ran\\
Đọc & không hỗ trợ & \verb|x = *p|, \verb|p->x| &
\verb|x = *p|, \verb|p->x| & \verb|x = *p|, \verb|p->x| &
\verb|x = *p|, \verb|p->x|, \verb|p[i]|\\
Ghi & \verb|*p = x| & không hỗ trợ & \verb|*p = x| &
\verb|*p = x| & \verb|*p = x|\\
Duyệt & \verb|++| & \verb|++| & \verb|++| & \verb|++|, \verb|--| &
\verb|++|, \verb|--|, \verb|+|, \verb|-|, \verb|+=|, \verb|-=|\\
So sánh & không hỗ trợ & \verb|== !=| & \verb|== !=| & \verb|== !=| &
\verb|== != < > <= >=|\\
\hline
\end{tabular}
\end{center}

Từ bảng này có thể thấy kiểu con trỏ chính là một kiểu bộ lặp ngẫu nhiên đặc
biệt. Với bộ lặp thông thường, các chức năng này được hiện thực bằng nạp
chồng toán tử.

Bộ lặp tiến tương thích với bộ lặp xuất và bộ lặp nhập; bộ lặp hai chiều
tương thích với bộ lặp tiến; bộ lặp ngẫu nhiên tương thích với bộ lặp hai
chiều. Quan hệ tương thích này có tính bắc cầu.

Kiểu cơ sở của bộ lặp được định nghĩa trong tập tin đầu \verb|iterator|. Vì
nó là một lớp trừu tượng, định nghĩa trông hơi lạ:

\begin{verbatim}
template <class Cat, class T, class Dist = ptrdiff_t,
    class Ptr = T*, class Ref = T&>
struct iterator {
    typedef Cat iterator_category;
    typedef T value_type;
    typedef Dist difference_type;
    typedef Ptr pointer;
    typedef Ref reference;
};
\end{verbatim}

Lớp \verb|iterator| chỉ cung cấp một giao diện cho kiểu bộ lặp; hiện thực cụ
thể sẽ được đặt ở những nơi cần đến.

Bộ lặp không chỉ dùng để truy cập phần tử. Trong toàn bộ STL, nó còn đóng vai
trò liên kết thuật toán và vật chứa. Vai trò này sẽ thấy rõ ở các mục sau.

Tập tin đầu \verb|utility| còn định nghĩa lớp khuôn mẫu \verb|pair| khá
thường dùng, để biểu diễn bộ đôi:

\begin{verbatim}
template <class T1, class T2> struct pair {
    typedef T1 first_type;
    typedef T2 second_type;
    T1 first;
    T2 second;
    pair() : first(T1()), second(T2()) { }
    pair(const T1& x, const T2& y) : first(x), second(y) { }
    template<class U, class V> pair(const pair<U, V>& p)
        : first(p.first), second(p.second) { }
};
template <class T1, class T2>
pair<T1, T2> make_pair(const T1& t1, const T2& t2) {
    return pair<T1, T2>(t1, t2);
}
\end{verbatim}

Lớp này sẽ xuất hiện nhiều lần khi mô tả thuật toán và vật chứa.

\subsection{Thuật toán}

Các thuật toán trong STL cung cấp một lượng lớn thao tác cơ bản, phần lớn là
thao tác trên dãy. Trước khi giới thiệu thuật toán STL, gọi tắt là thuật toán,
ta cần biết vài khái niệm nền tảng liên quan.

\subsubsection*{Dãy}

Đối tượng thao tác của hầu hết thuật toán là dãy. Trong STL, mô tả một dãy
cần hai bộ lặp: bộ lặp đầu trỏ tới phần tử đầu tiên của dãy, bộ lặp thứ hai
trỏ tới vị trí sau phần tử cuối cùng. Nói cách khác, khoảng nửa kín nửa mở do
hai bộ lặp làm đầu mút biểu diễn toàn bộ dãy.

Cách định nghĩa này ban đầu có thể lạ, nhưng nó đem lại nhiều tiện lợi. Với
dãy \([first,last)\), biểu thức \verb|last - first| biểu diễn độ dài dãy nếu
bộ lặp hỗ trợ phép trừ. Với mảng một chiều \verb|s| có \(n\) phần tử, ta dùng
\([s,s+n)\) để mô tả toàn bộ mảng. Khi mô tả vị trí chèn, bộ lặp \verb|p| có
nghĩa là chèn trước phần tử mà \verb|p| trỏ tới; như vậy \verb|first| là chèn
vào đầu dãy, \verb|last| là chèn vào cuối dãy, và mọi vị trí giữa cũng có bộ
lặp tương ứng.

Với các vật chứa STL ở mục sau, ta dùng các hàm thành viên do lớp vật chứa
cung cấp để tạo bộ lặp, rồi từ đó tạo dãy gồm toàn bộ hoặc một phần phần tử
trong vật chứa.

\subsubsection*{Đối tượng hàm}

Một số thuật toán ngoài việc mô tả dãy còn yêu cầu cung cấp đối tượng hàm.
Ví dụ:

\begin{verbatim}
template<class In, class Op>
    Op for_each(In first, In last, Op f);
\end{verbatim}

Hàm này thực hiện \verb|f(x)| một lần với mỗi phần tử \(x\) trong
\([first,last)\). Tất nhiên ta có thể dùng tên hàm làm tham số \verb|f|, nhưng
với một số thao tác, chẳng hạn tính tổng, có thể cần vài biến toàn cục. Nếu
định nghĩa các biến này toàn cục, ta phá vỡ nguyên tắc đóng gói. Cách giải
quyết là dùng đối tượng hàm. Với phép tính tổng, có thể viết lớp sau:

\begin{verbatim}
template<class T> class Sum {
    T res;
public:
    Sum(T i = T()) : res(i) { }
    void operator ()(const T& x) { res += x; }
    T result() const { return res; }
};
\end{verbatim}

Khi đó có thể tính tổng một dãy như sau:

\begin{verbatim}
Sum<int> sum;
for_each(s, s + n, sum(0));
cout << sum.result() << endl;
\end{verbatim}

Thực ra STL đã cung cấp thuật toán tính tổng; cách trên chỉ để minh họa cách
dùng đối tượng hàm.

Đối tượng hàm, kể cả hàm thường, có kiểu trả về là \verb|bool| được gọi là
vị từ, thường dùng để kiểm tra một quan hệ nào đó có thành lập hay không.

Trong tập tin đầu \verb|functional|, STL cung cấp một số đối tượng hàm cho
các phép toán thường dùng, đều là lớp khuôn mẫu.

\begin{center}
\begin{tabular}{lll}
\hline
\textbf{Tên lớp} & \textbf{Loại} & \textbf{Tác dụng}\\
\hline
\verb|equal_to| & hai ngôi & \verb|arg1 == arg2|\\
\verb|not_equal_to| & hai ngôi & \verb|arg1 != arg2|\\
\verb|greater| & hai ngôi & \verb|arg1 > arg2|\\
\verb|less| & hai ngôi & \verb|arg1 < arg2|\\
\verb|greater_equal| & hai ngôi & \verb|arg1 >= arg2|\\
\verb|less_equal| & hai ngôi & \verb|arg1 <= arg2|\\
\verb|logical_and| & hai ngôi & \verb|arg1 && arg2|\\
\verb|logical_or| & hai ngôi & \verb|arg1 || arg2|\\
\verb|logical_not| & một ngôi & \verb|!arg|\\
\verb|plus| & hai ngôi & \verb|arg1 + arg2|\\
\verb|minus| & hai ngôi & \verb|arg1 - arg2|\\
\verb|multiplies| & hai ngôi & \verb|arg1 * arg2|\\
\verb|divides| & hai ngôi & \verb|arg1 / arg2|\\
\verb|modulus| & hai ngôi & \verb|arg1 % arg2|\\
\verb|negate| & một ngôi & \verb|-arg|\\
\hline
\end{tabular}
\end{center}

\subsubsection*{Khái quát các thuật toán}

STL cung cấp hơn \(60\) thuật toán. Chúng đều là hàm khuôn mẫu và thường có
hai phiên bản nạp chồng: một phiên bản dùng toán tử mặc định, thường là bằng
hoặc nhỏ hơn tùy loại thuật toán; phiên bản còn lại cho phép tự chỉ định đối
tượng hàm thay cho toán tử mặc định.

Bảng sau liệt kê các thuật toán STL. Trừ bốn thuật toán cuối được định nghĩa
trong \verb|numeric|, các thuật toán còn lại nằm trong \verb|algorithm|.

\begin{longtable}{p{0.35\textwidth}p{0.57\textwidth}}
\hline
\textbf{Tên thuật toán} & \textbf{Tác dụng}\\
\hline
\endfirsthead
\hline
\textbf{Tên thuật toán} & \textbf{Tác dụng}\\
\hline
\endhead
\verb|for_each()| & Thực hiện một thao tác với mỗi phần tử của dãy\\
\verb|find()| & Tìm vị trí xuất hiện đầu tiên của một giá trị trong dãy\\
\verb|find_if()| & Tìm phần tử đầu tiên thỏa một vị từ\\
\verb|find_first_of()| & Tìm vị trí đầu tiên trong một dãy chứa phần tử bất kỳ của dãy khác\\
\verb|adjacent_find()| & Tìm cặp phần tử kề nhau đầu tiên bằng nhau\\
\verb|count()| & Đếm số lần một giá trị xuất hiện trong dãy\\
\verb|count_if()| & Đếm số phần tử thỏa một vị từ\\
\verb|mismatch()| & Tìm cặp phần tử khác nhau đầu tiên của hai dãy\\
\verb|equal()| & Kiểm tra hai dãy có hoàn toàn bằng nhau không\\
\verb|search()| & Tìm vị trí xuất hiện đầu tiên của một dãy con\\
\verb|find_end()| & Tìm vị trí xuất hiện cuối cùng của một dãy con\\
\verb|search_n()| & Tìm vị trí một giá trị xuất hiện liên tiếp \(n\) lần\\
\verb|transform()| & Thực hiện thao tác trên từng phần tử và tạo dãy xuất tương ứng\\
\verb|copy()| & Sao chép dãy sang dãy khác từ trước ra sau\\
\verb|copy_backward()| & Sao chép dãy sang dãy khác từ sau ra trước\\
\verb|swap()| & Đổi giá trị hai phần tử\\
\verb|iter_swap()| & Đổi giá trị hai phần tử mà hai bộ lặp trỏ tới\\
\verb|swap_ranges()| & Đổi phần tử giữa hai dãy\\
\verb|replace()| & Thay một giá trị trong dãy bằng giá trị khác\\
\verb|replace_if()| & Thay các phần tử thỏa một vị từ bằng giá trị khác\\
\verb|replace_copy()| & Như \verb|replace()| nhưng ghi kết quả sang dãy khác\\
\verb|replace_copy_if()| & Như \verb|replace_if()| nhưng ghi kết quả sang dãy khác\\
\verb|fill()| & Gán mọi phần tử trong dãy bằng một giá trị\\
\verb|fill_n()| & Như \verb|fill()| nhưng dãy được mô tả bằng bộ lặp đầu và độ dài\\
\verb|generate()| & Gán mọi phần tử bằng giá trị trả về của bộ sinh\\
\verb|generate_n()| & Như \verb|generate()| nhưng dãy được mô tả bằng bộ lặp đầu và độ dài\\
\verb|remove()| & Loại khỏi dãy các phần tử bằng một giá trị\\
\verb|remove_if()| & Loại khỏi dãy các phần tử thỏa một vị từ\\
\verb|remove_copy()| & Như \verb|remove()| nhưng ghi kết quả sang dãy khác\\
\verb|remove_copy_if()| & Như \verb|remove_if()| nhưng ghi kết quả sang dãy khác\\
\verb|unique()| & Loại các phần tử bằng nhau kề nhau bị thừa\\
\verb|unique_copy()| & Như \verb|unique()| nhưng ghi kết quả sang dãy khác\\
\verb|reverse()| & Đảo ngược dãy\\
\verb|reverse_copy()| & Như \verb|reverse()| nhưng ghi kết quả sang dãy khác\\
\verb|rotate()| & Xoay dãy\\
\verb|rotate_copy()| & Như \verb|rotate()| nhưng ghi kết quả sang dãy khác\\
\verb|random_shuffle()| & Xáo trộn ngẫu nhiên các phần tử trong dãy\\
\verb|sort()| & Sắp xếp các phần tử trong dãy\\
\verb|stable_sort()| & Như \verb|sort()| nhưng bảo đảm tính ổn định\\
\verb|partial_sort()| & Sắp xếp phần đầu của dãy\\
\verb|partial_sort_copy()| & Như \verb|partial_sort()| nhưng ghi kết quả sang dãy khác\\
\verb|nth_element()| & Đưa phần tử nhỏ thứ \(n\) về đúng vị trí của nó sau khi sắp xếp\\
\verb|lower_bound()| & Tìm nhị phân trong dãy có thứ tự, trả về vị trí xuất hiện đầu tiên\\
\verb|upper_bound()| & Tìm nhị phân trong dãy có thứ tự, trả về vị trí sau lần xuất hiện cuối\\
\verb|equal_range()| & Tìm nhị phân và trả về dãy con gồm mọi phần tử bằng giá trị cần tìm\\
\verb|binary_search()| & Kiểm tra một giá trị có xuất hiện trong dãy có thứ tự không\\
\verb|merge()| & Trộn hai dãy có thứ tự\\
\verb|inplace_merge()| & Trộn hai dãy có thứ tự liên tiếp\\
\verb|partition()| & Sắp lại dãy để phần tử thỏa vị từ đứng trước\\
\verb|stable_partition()| & Như \verb|partition()| nhưng bảo đảm tính ổn định\\
\verb|includes()| & Kiểm tra một dãy có là dãy con, không cần liên tiếp, của dãy khác không\\
\verb|set_union()| & Hợp hai dãy có thứ tự theo nghĩa tập hợp\\
\verb|set_intersection()| & Giao hai dãy có thứ tự theo nghĩa tập hợp\\
\verb|set_difference()| & Hiệu hai dãy có thứ tự theo nghĩa tập hợp\\
\verb|set_symmetric_difference()| & Hiệu đối xứng hai dãy có thứ tự\\
\verb|make_heap()| & Điều chỉnh một dãy thành heap\\
\verb|push_heap()| & Chèn một phần tử vào heap\\
\verb|pop_heap()| & Lấy phần tử đỉnh heap ra\\
\verb|sort_heap()| & Điều chỉnh heap thành dãy có thứ tự\\
\verb|min()| & Trả về giá trị nhỏ hơn trong hai giá trị\\
\verb|max()| & Trả về giá trị lớn hơn trong hai giá trị\\
\verb|min_element()| & Trả về vị trí phần tử nhỏ nhất trong dãy\\
\verb|max_element()| & Trả về vị trí phần tử lớn nhất trong dãy\\
\verb|lexicographical_compare()| & So sánh hai dãy theo thứ tự từ điển\\
\verb|next_permutation()| & Trả về hoán vị kế tiếp theo thứ tự từ điển\\
\verb|prev_permutation()| & Trả về hoán vị trước theo thứ tự từ điển\\
\verb|accumulate()| & Tính tổng các phần tử của dãy\\
\verb|inner_product()| & Tính tích vô hướng của hai dãy\\
\verb|partial_sum()| & Tính dãy tổng từng phần\\
\verb|adjacent_difference()| & Tính dãy hiệu của các phần tử kề nhau\\
\hline
\end{longtable}

Tính ổn định trong sắp xếp nghĩa là trong dãy sau khi sắp xếp, các phần tử có
khóa bằng nhau giữ nguyên thứ tự tương đối như trong dãy ban đầu.

Bảng trên chỉ giới thiệu chức năng cơ bản. Sau đây là vài thuật toán thường
dùng. Trong thi tin học, thuật toán được dùng nhiều nhất là sắp xếp. Hàm
\verb|sort()| của STL làm việc này rất thuận tiện.

\begin{verbatim}
template<class Ran> void sort(Ran first, Ran last);
template<class Ran, class Cmp>
    void sort(Ran first, Ran last, Cmp cmp);
\end{verbatim}

Với kiểu đã định nghĩa toán tử nhỏ hơn, chỉ cần viết \verb|sort(s, s + n)| là
sắp xếp xong. Nếu chưa định nghĩa toán tử nhỏ hơn hoặc muốn dùng quan hệ khác
để sắp xếp, dùng phiên bản nạp chồng thứ hai. Khi đó ngoài dãy cần sắp xếp,
ta đưa thêm một vị từ làm căn cứ so sánh. Ví dụ
\verb|sort(s, s + n, greater<int>())| sẽ sắp xếp giảm dần.

\verb|sort()| thường được hiện thực bằng quicksort, thời gian trung bình
\(O(n\log n)\). Trường hợp xấu nhất là \(O(n^2)\), nhưng xác suất gặp thường
rất nhỏ.

Tìm nhị phân cũng là thuật toán rất thường dùng trong thi tin học. STL có bốn
thuật toán liên quan:

\begin{verbatim}
template<class For, class T> bool
    binary_search(For first, For last, const T& val);
template<class For, class T, class Cmp> bool
    binary_search(For first, For last, const T& val, Cmp cmp);
template<class For, class T> For
    lower_bound(For first, For last, const T& val);
template<class For, class T, class Cmp> For
    lower_bound(For first, For last, const T& val, Cmp cmp);
template<class For, class T> For
    upper_bound(For first, For last, const T& val);
template<class For, class T, class Cmp> For
    upper_bound(For first, For last, const T& val, Cmp cmp);
template<class For, class T> pair<For, For>
    equal_range(For first, For last, const T& val);
template<class For, class T, class Cmp> pair<For, For>
    equal_range(For first, For last, const T& val, Cmp cmp);
\end{verbatim}

Nhìn tên thì thuật toán tìm nhị phân phải là \verb|binary_search()|, nhưng
thực tế hàm này không quá hữu dụng vì chỉ trả về một giá trị có xuất hiện
trong dãy hay không. Ba hàm sau mới thường hữu ích hơn.

Nếu giá trị cần tìm xuất hiện trong dãy, \verb|lower_bound()| trả về vị trí
phần tử đầu tiên bằng giá trị đó, còn \verb|upper_bound()| trả về vị trí ngay
sau phần tử cuối cùng bằng giá trị đó. Cách định nghĩa này phù hợp với định
nghĩa dãy trong STL. Gọi giá trị trả về của \verb|lower_bound()| là \verb|l|
và của \verb|upper_bound()| là \verb|u|, thì dãy \([l,u)\) chính là dãy con
gồm mọi phần tử bằng giá trị cần tìm. Nếu giá trị cần tìm không xuất hiện,
hai hàm trả về cùng một vị trí, sao cho nếu chèn giá trị cần tìm vào đó thì
dãy vẫn giữ thứ tự. \verb|equal_range()| là phiên bản gộp của
\verb|lower_bound()| và \verb|upper_bound()|. Cách dùng chi tiết các thuật
toán khác nên xem tài liệu.

\subsubsection*{Ví dụ: dãy con tăng dài nhất}

\paragraph{Mô tả bài toán.}
Cho dãy số nguyên độ dài \(n\),
\(A=\{a_1,a_2,\ldots,a_n\}\). Hãy tìm số nguyên lớn nhất \(m\) sao cho tồn
tại một dãy chỉ số \(P=\{p_1,p_2,\ldots,p_m\}\), thỏa
\(1\le p_1<p_2<\cdots<p_m\le n\) và
\(a_{p_1}<a_{p_2}<\cdots<a_{p_m}\).

\paragraph{Dữ liệu vào.}
Tập tin \verb|LMIS.in|. Dòng thứ nhất là số nguyên dương \(n\le 30000\), độ
dài dãy ban đầu. Dòng thứ hai là \(n\) số nguyên dương \(a_i\le 1000000000\).

\paragraph{Dữ liệu ra.}
Tập tin \verb|LMIS.out|. Một dòng chứa số nguyên \(m\), độ dài của dãy con
tăng dài nhất.

\begin{verbatim}
Input
10
3 1 4 1 5 9 2 6 5 3

Output
4
\end{verbatim}

Đây là một bài quy hoạch động điển hình. Gọi \(f_i\) là độ dài dãy con tăng
dài nhất kết thúc tại phần tử thứ \(i\) của dãy ban đầu. Để tiện, đặt
\(a_0=-\infty\), \(f_0=0\). Phương trình chuyển trạng thái:
\[
f_i=\max_{0\le j<i,\ a_j<a_i}\{f_j+1\}.
\]
Kết quả cần tìm là giá trị lớn nhất trong các \(f_i\). Chương trình hiện thực
thuật toán này như sau; vì chỉ số mảng \cpp{} bắt đầu từ \(0\), chỉ số được
điều chỉnh đôi chút.

\begin{verbatim}
//: C03:LMIS1.cpp
// Longest Monotonically Increasing Subsequence
// Original Algorithm O(n^2)
#include <fstream>
#include <algorithm>
using namespace std;
ifstream fin("LMIS.in");
ofstream fout("LMIS.out");
const int max_n = 30000;
int n;
int a[max_n];
int f[max_n];
int main() {
    fin >> n;
    for(int i = 0; i < n; i++)
        fin >> a[i];
    for(int i = 0; i < n; i++) {
        f[i] = 1;
        for(int j = 0; j < i; j++)
            if(a[j] < a[i])
                f[i] = max(f[i], f[j] + 1);
    }
    fout << *max_element(f, f + n) << endl;
} ///:~
\end{verbatim}

Thuật toán này có thời gian \(O(n^2)\), không đáp ứng dữ liệu. Để cải tiến,
ta dùng một mảng phụ. Gọi \(g_i\) là giá trị nhỏ nhất của phần tử cuối trong
tất cả dãy con tăng độ dài \(i\) đã thấy đến hiện tại. Dễ biết
\(g_{i-1}\le g_i\) với \(1\le i\le n\), có thể chứng minh bằng phản chứng;
tức \(\{g_i\}\) là dãy không giảm. Khi đã biết \(\{g_i\}\) đến phần tử thứ
\(i-1\), \(f_i\) bằng vị trí của phần tử đầu tiên trong \(\{g_i\}\) lớn hơn
hoặc bằng \(a_i\), cũng chính là vị trí do \verb|lower_bound()| trả về. Sau
đó cập nhật \(g_{f_i}=a_i\). Ban đầu đặt \(g_i=\infty\) với \(1\le i\le n\).
Mỗi lần tìm vị trí mất \(O(\log n)\), cập nhật mất \(O(1)\), nên tổng thời
gian là \(O(n\log n)\).

\begin{verbatim}
//: C03:LMIS2.cpp
// Longest Monotonically Increasing Subsequence
// Improved Algorithm O(n log n)
#include <fstream>
#include <algorithm>
using namespace std;
ifstream fin("LMIS.in");
ofstream fout("LMIS.out");
const int infinity = 2000000000;
const int max_n = 30000;
int n;
int a[max_n];
int f[max_n];
int g[max_n];
int main() {
    fin >> n;
    for(int i = 0; i < n; i++)
        fin >> a[i];
    fill(g, g + n, infinity);
    for(int i = 0; i < n; i++) {
        int j = lower_bound(g, g + n, a[i]) - g;
        f[i] = j + 1;
        g[j] = a[i];
    }
    fout << *max_element(f, f + n) << endl;
} ///:~
\end{verbatim}

So sánh hai chương trình cho thấy nhờ STL, chương trình cải tiến thậm chí còn
đơn giản hơn chương trình ban đầu, thể hiện rõ ưu điểm giảm độ phức tạp khi
lập trình.

\subsection{Vật chứa}

Thuật toán chỉ phát huy tác dụng khi tác động lên cấu trúc dữ liệu. Ngoài
nhiều thuật toán, STL còn cung cấp một số cấu trúc dữ liệu thường dùng. Vì
chúng đều dùng để chứa một nhóm đối tượng cùng loại, chúng được gọi là vật
chứa. Bảng sau liệt kê các vật chứa STL phổ biến.

\begin{center}
\small
\begin{tabular}{lllll}
\hline
\textbf{Tên} & \textbf{Mô tả} & \textbf{Loại} & \textbf{Tập tin đầu} &
\textbf{Bộ lặp}\\
\hline
\verb|vector| & vector & vật chứa thường & \verb|vector| & ngẫu nhiên\\
\verb|deque| & hàng đợi hai đầu & vật chứa thường & \verb|deque| & ngẫu nhiên\\
\verb|list| & danh sách liên kết & vật chứa thường & \verb|list| & hai chiều\\
\verb|stack| & ngăn xếp & bộ thích nghi & \verb|stack| & không cung cấp\\
\verb|queue| & hàng đợi & bộ thích nghi & \verb|queue| & không cung cấp\\
\verb|priority_queue| & hàng đợi ưu tiên & bộ thích nghi & \verb|queue| & không cung cấp\\
\verb|set| & tập hợp & vật chứa kết hợp & \verb|set| & hai chiều\\
\verb|multiset| & đa tập & vật chứa kết hợp & \verb|set| & hai chiều\\
\verb|map| & ánh xạ & vật chứa kết hợp & \verb|map| & hai chiều\\
\verb|multimap| & đa ánh xạ & vật chứa kết hợp & \verb|map| & hai chiều\\
\hline
\end{tabular}
\end{center}

Bộ thích nghi vật chứa chỉ cung cấp một giao diện hàm, không cung cấp bộ lặp,
và bên trong chứa một vật chứa thường để thực hiện lưu trữ phần tử. Vật chứa
kết hợp thường có phần tử gồm khóa và giá trị được ánh xạ, cho phép tìm nhanh
theo khóa.

Mọi vật chứa đều được hiện thực bằng lớp khuôn mẫu, nên khi dùng phải chỉ
định kiểu phần tử. Điều này có thể làm tên kiểu dài và viết bất tiện; khi cần
có thể dùng \verb|typedef| để đơn giản hóa.

\begin{verbatim}
typedef vector<string> str_vec;
\end{verbatim}

Vật chứa STL tự động cấp phát không gian cho phần tử. Ví dụ \verb|vector| sẽ
giữ trước một phần không gian; khi chèn quá nhiều phần tử, nó tự cấp phát một
vùng lớn hơn và sao chép phần tử cũ sang. Khi đối tượng vật chứa bị hủy, không
gian của phần tử cũng tự được thu hồi.

Với các vật chứa này, STL cung cấp hệ thống kiểu thành viên và giao diện hàm
cơ bản tương đối giống nhau, dù từng vật chứa có khác biệt.

\begin{longtable}{p{0.30\textwidth}p{0.18\textwidth}p{0.43\textwidth}}
\hline
\textbf{Tên} & \textbf{Loại} & \textbf{Mô tả ngắn}\\
\hline
\endfirsthead
\hline
\textbf{Tên} & \textbf{Loại} & \textbf{Mô tả ngắn}\\
\hline
\endhead
\verb|value_type| & kiểu thành viên & Kiểu phần tử\\
\verb|iterator| & kiểu thành viên & Kiểu bộ lặp, giống \verb|value_type*|\\
\verb|const_iterator| & kiểu thành viên & Kiểu bộ lặp hằng, giống \verb|const value_type*|\\
\verb|reverse_iterator| & kiểu thành viên & Kiểu bộ lặp ngược\\
\verb|const_reverse_iterator| & kiểu thành viên & Kiểu bộ lặp ngược hằng\\
\verb|reference| & kiểu thành viên & Tham chiếu phần tử, bằng \verb|value_type&|\\
\verb|const_reference| & kiểu thành viên & Tham chiếu hằng phần tử\\
\verb|key_type| & kiểu thành viên & Kiểu khóa, chỉ dùng cho vật chứa kết hợp\\
\verb|mapped_type| & kiểu thành viên & Kiểu được ánh xạ, chỉ dùng cho vật chứa kết hợp\\
\verb|begin()| & thao tác bộ lặp & Trả về bộ lặp trỏ tới phần tử đầu\\
\verb|end()| & thao tác bộ lặp & Trả về bộ lặp trỏ tới vị trí sau phần tử cuối\\
\verb|rbegin()| & thao tác bộ lặp & Trả về bộ lặp ngược trỏ tới phần tử cuối\\
\verb|rend()| & thao tác bộ lặp & Trả về bộ lặp ngược trỏ tới trước phần tử đầu\\
\verb|front()| & thao tác phần tử & Phần tử đầu\\
\verb|back()| & thao tác phần tử & Phần tử cuối\\
\verb|operator []()| & thao tác phần tử & Toán tử chỉ số\\
\verb|push_back()| & ngăn xếp/hàng đợi & Thêm một phần tử vào cuối\\
\verb|pop_back()| & ngăn xếp/hàng đợi & Xóa phần tử cuối\\
\verb|push_front()| & ngăn xếp/hàng đợi & Thêm một phần tử vào đầu, \verb|vector| không hỗ trợ\\
\verb|pop_front()| & ngăn xếp/hàng đợi & Xóa phần tử đầu, \verb|vector| không hỗ trợ\\
\verb|insert()| & thao tác danh sách & Chèn một số phần tử\\
\verb|erase()| & thao tác danh sách & Xóa một số phần tử\\
\verb|clear()| & thao tác danh sách & Xóa sạch vật chứa\\
\verb|size()| & thao tác khác & Trả về số phần tử\\
\verb|empty()| & thao tác khác & Kiểm tra vật chứa rỗng\\
\verb|resize()| & thao tác khác & Điều chỉnh số phần tử\\
\verb|find()| & vật chứa kết hợp & Tìm một khóa\\
\verb|lower_bound()| & vật chứa kết hợp & Tìm khóa, trả về vị trí đầu phù hợp\\
\verb|upper_bound()| & vật chứa kết hợp & Tìm khóa, trả về vị trí sau vị trí cuối phù hợp\\
\verb|equal_range()| & vật chứa kết hợp & Tìm khóa, trả về khoảng phù hợp\\
\hline
\end{longtable}

Sau đây là vài vật chứa thường dùng.

\verb|vector| là mảng một chiều thông minh có thể tự điều chỉnh không gian
lưu trữ. \verb|vector| còn tối ưu cách lưu \verb|vector<bool>|, dùng mỗi bit
để lưu một giá trị \verb|bool|. Tuy nhiên trong thi tin học, quy mô dữ liệu
thường đã biết trước, hơn nữa việc điều chỉnh không gian của \verb|vector|
phải di chuyển dữ liệu và tốn khá nhiều thời gian, nên \verb|vector| không
phải lúc nào cũng thực dụng nhất.

\verb|deque| được hiện thực bằng mảng con trỏ và chừa một lượng không gian ở
cả hai đầu, nên thao tác chèn/xóa ở đầu và cuối rất nhanh. \verb|deque| thường
dùng làm vật chứa bên trong cho các bộ thích nghi như \verb|stack| và
\verb|queue|.

\verb|list| thường được hiện thực bằng danh sách liên kết hai chiều, giúp ta
tránh phải viết nhiều thao tác danh sách cơ bản lặp lại và giảm độ phức tạp
khi lập trình.

\verb|stack|, \verb|queue| và \verb|priority_queue| đều là bộ thích nghi vật
chứa. Vì giao diện ít hàm, công dụng của \verb|stack| và \verb|queue| có phần
bị giới hạn, nhưng đủ cho các tình huống thông thường. \verb|priority_queue|
là hàng đợi ưu tiên hiện thực bằng heap nhị phân. Khác hàng đợi thường, khi
lấy phần tử, hàng đợi ưu tiên luôn lấy phần tử có độ ưu tiên cao nhất trước;
với cách so sánh mặc định của \cpp{} chuẩn, phần tử lớn nhất đứng đầu.
\verb|priority_queue| rất hữu ích trong bài mô phỏng, vì mấu chốt của mô
phỏng là xử lý đúng thứ tự sự kiện. Ta có thể xem các sự kiện trong đề là đối
tượng, gán độ ưu tiên theo quy tắc đề, nạp chồng toán tử so sánh, rồi lần
lượt đẩy vào hàng đợi ưu tiên và lấy ra xử lý. Tuy nhiên do giới hạn giao
diện, hàng đợi ưu tiên khó xử lý tình huống phần tử trong hàng đợi cần thay
đổi, chẳng hạn trong thuật toán Dijkstra hoặc Prim.

Các vật chứa kết hợp của STL rất mạnh. Bốn vật chứa này đều được hiện thực
bằng cây nhị phân cân bằng, thường là cây đỏ-đen. Phần tử trong vật chứa có
thứ tự, và các thao tác chèn, xóa, tìm đều hoàn thành trong \(O(\log n)\).
Trong \verb|set| và \verb|multiset|, bản thân phần tử trực tiếp tham gia sắp
xếp, hay nói cách khác phần tử chính là khóa. Trong \verb|map| và
\verb|multimap|, chỉ khóa tham gia sắp xếp; giá trị được ánh xạ đi kèm như dữ
liệu phụ và có thể sửa. Ngoài ra, \verb|set| và \verb|map| không cho phép
khóa bằng nhau, còn \verb|multiset| và \verb|multimap| cho phép khóa bằng
nhau.

Các vật chứa kết hợp đều có thể tìm phần tử theo khóa. Các hàm thành viên tìm
kiếm gồm \verb|find()|, \verb|lower_bound()|, \verb|upper_bound()| và
\verb|equal_range()|. Khi không tìm thấy khóa yêu cầu, \verb|find()| trả về
giá trị của \verb|end()| để biểu thị thất bại; ba hàm còn lại có giá trị trả
về giống các hàm cùng tên trong \verb|algorithm|. Vì trong \verb|map| khóa và
giá trị được ánh xạ tương ứng một-một, lớp \verb|map| nạp chồng toán tử chỉ
số, làm việc truy cập phần tử thuận tiện hơn.

\begin{verbatim}
map<string, int> week;
week["Sunday"] = 0;
week["Monday"] = 1;
week["Tuesday"] = 2;
week["Wednesday"] = 3;
week["Thursday"] = 4;
week["Friday"] = 5;
week["Saturday"] = 6;
cout << week["Thursday"] << endl;
\end{verbatim}

Khi dùng toán tử chỉ số của \verb|map|, nếu phần tử ứng với khóa trong ngoặc
vuông tồn tại, tham chiếu đến giá trị được ánh xạ của phần tử đó được trả về.
Nếu phần tử ứng với khóa đó chưa tồn tại, một phần tử mới có khóa ấy được
chèn vào, giá trị được ánh xạ được khởi tạo bằng giá trị mặc định của kiểu,
và tham chiếu đến giá trị được ánh xạ được trả về.

\subsubsection*{Ví dụ: thanh toán hóa đơn}

\paragraph{Mô tả bài toán.}
Bill gần đây gặp một việc phiền. Mỗi buổi sáng anh nhận một số hóa đơn, có
thể không nhận hóa đơn nào; mỗi hóa đơn có một mệnh giá. Buổi chiều, anh chọn
hai hóa đơn chưa thanh toán có mệnh giá lớn nhất và nhỏ nhất rồi trả chúng.
Các hóa đơn chưa trả được giữ lại đến ngày sau. Biết số lượng và mệnh giá hóa
đơn Bill nhận mỗi ngày, hãy giúp anh đưa ra thứ tự thanh toán.

\paragraph{Dữ liệu vào.}
Dòng thứ nhất là số nguyên dương \(N\le 30000\), tổng số ngày. Từ dòng thứ
hai đến dòng \(N+1\), mỗi dòng mô tả một ngày: trước hết là số nguyên không âm
\(M\), số hóa đơn nhận trong ngày, sau đó là \(M\) số nguyên dương nhỏ hơn
\(1000000000\), mệnh giá các hóa đơn. Dữ liệu bảo đảm mỗi ngày đều có thể
thanh toán hai hóa đơn, và đến ngày cuối mọi hóa đơn đều được thanh toán.

\paragraph{Dữ liệu ra.}
Gồm \(N\) dòng, mỗi dòng hai số nguyên cách nhau bởi khoảng trắng, lần lượt là
mệnh giá nhỏ nhất và lớn nhất được thanh toán trong ngày.

\begin{verbatim}
Input
4
3 3 6 5
2 8 2
3 7 1 7
0

Output
3 6
2 8
1 7
5 7
\end{verbatim}

Từ giới hạn dữ liệu có thể thấy thuật toán cần tốt hơn \(O(n^2)\), và
\(O(n\log n)\) chấp nhận được. Có nhiều thuật toán \(O(n\log n)\). Chẳng hạn
có thể tạo hai heap, một heap lớn nhất và một heap nhỏ nhất, rồi thiết lập ánh
xạ giữa các phần tử tương ứng. Khi nhận hóa đơn, đẩy mệnh giá vào cả hai heap
và điều chỉnh vị trí; khi thanh toán, lấy phần tử đầu của mỗi heap và điều
chỉnh. Vì vị trí của phần tử trong heap cần thay đổi sau khi đưa vào, khó
dùng \verb|priority_queue| của STL để hiện thực; nếu chọn thuật toán này,
phải tự viết mã heap.

Một thuật toán \(O(n\log n)\) khác là dùng cây nhị phân cân bằng. Tự viết cây
cân bằng còn phức tạp hơn cách trên. Nhưng tập hợp của STL được hiện thực
bằng cây đỏ-đen, nên nhờ nó chương trình rất đơn giản. Vì mệnh giá hóa đơn có
thể trùng nhau, phải dùng \verb|multiset|.

\begin{verbatim}
//: C03:Bills.cpp
// Pay 2 Bills Every Day
#include <iostream>
#include <set>
using namespace std;
int main() {
    int n;
    cin >> n;
    multiset<int> bills;
    while(n--) {
        int m;
        cin >> m;
        while(m--) {
            int a;
            cin >> a;
            bills.insert(a);
        }
        cout << *bills.begin() << ' ' << *--bills.end() << endl;
        bills.erase(bills.begin());
        bills.erase(--bills.end());
    }
} ///:~
\end{verbatim}

Ví dụ này một lần nữa thể hiện sức mạnh của STL trong việc giảm độ phức tạp
khi lập trình.

Ngoài các vật chứa STL nói trên, \cpp{} chuẩn còn cung cấp \verb|bitset|, tập
bit tương tự \verb|set| của \pascal{}, và \verb|valarray|, mảng số học có thể
thực hiện phép toán trên một nhóm dữ liệu cùng kiểu và có thể dùng để hiện
thực ma trận. Nhiều trình biên dịch \cpp{} dùng bản STL do SGI hiện thực, còn
cung cấp \verb|slist|, danh sách liên kết đơn tiết kiệm không gian,
\verb|hash_set|, \verb|hash_multiset|, \verb|hash_map|, \verb|hash_multimap|,
các tập hợp và ánh xạ hiện thực bằng bảng băm, cùng \verb|rope|, một cấu trúc
dữ liệu chuỗi hỗ trợ chèn và xóa trong thời gian logarit. Các vật chứa này
đều thực dụng; độc giả quan tâm có thể xem thêm \url{http://www.sgi.com/tech/stl}.

\subsection{Kết chương}

Ba mục trước đã trình bày ba thành phần quan trọng của STL: bộ lặp, thuật
toán và vật chứa. Nhưng ba phần này không độc lập; chỉ khi kết hợp hữu cơ với
nhau, STL mới phát huy đúng tác dụng. Khi vận dụng tổng hợp, cần chú ý một số
điểm.

Trước hết, ưu tiên dùng thuật toán do chính vật chứa cung cấp. Một số lớp vật
chứa có các hàm thành viên với chức năng gần giống thuật toán tổng quát, và
nên ưu tiên dùng chúng. Thuật toán tổng quát được thiết kế cho dãy nói chung,
phải cân nhắc tính phổ quát nên khó tối ưu sâu; còn hàm thành viên của lớp
vật chứa có thể dựa trên đặc điểm riêng của vật chứa để tối ưu hoặc làm những
việc thuật toán tổng quát không làm được. Ví dụ lớp \verb|list| có các hàm
thành viên \verb|remove()|, \verb|unique()|, \verb|merge()|,
\verb|reverse()|, \verb|sort()|, nên ưu tiên dùng. Thuật toán tổng quát
\verb|sort()| không thể sắp xếp phần tử trong \verb|list| vì \verb|list| chỉ
cung cấp bộ lặp hai chiều, còn \verb|sort()| cần bộ lặp ngẫu nhiên.

Thứ hai, khi kết hợp vật chứa và thuật toán cần chú ý dãy xuất. Nhiều thuật
toán yêu cầu chỉ định dãy xuất, thường biểu diễn bằng bộ lặp trỏ tới đầu dãy
xuất. Trong thuật toán, bộ lặp này liên tục được tăng để xuất toàn bộ dãy. Vì
bộ lặp thông thường không kiểm tra vượt biên, nếu không bảo đảm vật chứa đã có
đủ phần tử để lưu dãy xuất, không nên dùng trực tiếp bộ lặp của vật chứa để
chỉ định dãy xuất. Cách đúng là dùng bộ lặp chèn trong \verb|iterator|.

Có ba loại bộ lặp chèn: bộ lặp chèn đầu, bộ lặp chèn cuối và bộ lặp chèn
thông thường. Với vật chứa \verb|c|, chúng lần lượt được tạo bởi
\verb|front_inserter(c)|, \verb|back_inserter(c)| và \verb|inserter(c, i)|,
trong đó \verb|i| là một bộ lặp của \verb|c|. Chúng đều là bộ lặp xuất, và
khi tăng sẽ thực hiện thao tác chèn tương ứng: \verb|push_front()|,
\verb|push_back()| hoặc \verb|insert()|.

\begin{verbatim}
unique_copy(v.begin(), v.end(), back_inserter(v2));
set_union(s1.begin(), s1.end(), s2.begin(), s2.end(),
    inserter(s, s.begin()));
\end{verbatim}

Câu lệnh đầu xóa các phần tử lặp kề nhau trong vector \verb|v| và thêm kết
quả vào cuối vector \verb|v2|. Câu lệnh thứ hai lấy hợp của hai tập hợp
\verb|s1| và \verb|s2| rồi thêm kết quả vào tập hợp \verb|s|; ở đây
\verb|s.begin()| chỉ để đủ số tham số, vì phần tử trong \verb|set| vẫn được
sắp theo thứ tự, không phụ thuộc vị trí chèn.

Qua chương này ta đã có hiểu biết nhất định về STL trong \cpp{}. STL có nhiều
ưu điểm: giảm độ phức tạp khi lập trình và nâng cao xác suất viết đúng mã.
Tuy nhiên cái gì có lợi cũng có mặt trái. Vì STL dùng cơ chế lập trình tổng
quát như lớp khuôn mẫu và hàm khuôn mẫu, chỉ cần chương trình có lỗi nhỏ cũng
có thể sinh ra rất nhiều thông báo lỗi kỳ lạ lúc biên dịch, và vị trí báo lỗi
thường nằm trong tập tin đầu của STL, vì lỗi khuôn mẫu thường chỉ lộ ra khi
đặc biệt hóa. Lại vì STL dùng đóng gói lớp, việc gỡ lỗi động cũng bất tiện
hơn, do ta khó xem trực tiếp các phần tử bên trong vật chứa.

Một số đề nghị khi dùng STL trong thi tin học:
\begin{enumerate}
  \item Dùng nhiều thuật toán STL để giảm độ phức tạp khi lập trình và tăng
  độ đúng đắn.
  \item Trong điều kiện tương đương, ưu tiên dùng mảng dựng sẵn của \cpp{}
  thay vì vật chứa STL, để tiện quan sát khi gỡ lỗi.
  \item Cố gắng dùng kiểm tra tĩnh nhiều hơn kiểm tra động.
  \item Khi gỡ lỗi động, dùng cách xuất ra màn hình để xem các nội dung mà
  trình gỡ lỗi không xem được.
\end{enumerate}

\section*{Tổng kết}
\addcontentsline{toc}{section}{Tổng kết}

Quá trình học ngôn ngữ lập trình có nhiều điểm rất giống quá trình học ngôn
ngữ tự nhiên.

Chương thứ nhất trình bày quá trình chuyển từ \pascal{} sang \cpp{}. Khi học,
ta phải tận dụng đầy đủ các điểm giống nhau giữa hai ngôn ngữ và nhận rõ các
điểm khác biệt, nhờ đó tăng tốc độ chuyển đổi. Điều này giống như học ngoại
ngữ trên nền tiếng mẹ đẻ: tận dụng tốt điểm giống giữa các ngôn ngữ làm việc
học ngoại ngữ đơn giản hơn; nếu không xử lý tốt các điểm khác, việc học sẽ
gặp khó khăn. Muốn học ngoại ngữ tốt cần luyện tập nhiều; tương tự, chỉ khi
liên tục dùng một ngôn ngữ lập trình để viết chương trình, ta mới thật sự
nâng cao mức độ nắm vững ngôn ngữ đó.

Chương thứ hai và thứ ba đi sâu hơn, trình bày một số nội dung nâng cao của
\cpp{}. Điều này giống như những tầng văn hóa riêng trong ngoại ngữ; chỉ khi
hiểu những tầng ấy, ta mới thật sự nắm một ngôn ngữ. Ngôn ngữ lập trình cũng
vậy. Nội dung chương thứ hai và thứ ba giúp ta dùng \cpp{} tốt hơn, đặc biệt
là STL trong chương thứ ba, vốn giảm đáng kể độ phức tạp khi lập trình, nâng
cao độ đúng của mã và rất hữu ích cho thi tin học.

Lập trình trong thi tin học lại giống viết văn. Nền tảng của viết văn là nắm
tốt một ngôn ngữ tự nhiên; nền tảng của lập trình là nắm tốt một ngôn ngữ lập
trình. Nhưng muốn viết bài văn hay thì chỉ nắm ngôn ngữ tự nhiên là chưa đủ,
còn cần ý tưởng tinh tế và cách diễn đạt đẹp. Lập trình trong thi tin học
cũng vậy: nắm vững ngôn ngữ lập trình rất quan trọng, nhưng cốt lõi của lập
trình là thiết kế thuật toán. Chỉ khi kết hợp hữu cơ hai mặt này, ta mới có
thể vượt qua các bài toán trong thi tin học và đạt kết quả tốt.

\section*{Tài liệu tham khảo}
\addcontentsline{toc}{section}{Tài liệu tham khảo}

\begin{enumerate}
  \item \emph{The C++ Programming Language}, Special Edition.
  \item \emph{Thinking in C++}, 2nd Edition.
  \item \emph{Introduction to Algorithms}, 2nd Edition.
\end{enumerate}

\end{document}
